\documentclass[11pt]{article}
\usepackage[T1]{fontenc}
\usepackage{lmodern}
\usepackage[margin=1.05in]{geometry}
\usepackage{amsmath,amssymb,amsthm,mathtools,mathrsfs}
\usepackage{microtype}
\usepackage{enumitem}
\usepackage[colorlinks=true,linkcolor=blue,citecolor=blue,urlcolor=blue]{hyperref}
\usepackage{bookmark}
\hypersetup{pdftitle={Lech's multiplicity conjecture},pdfauthor={OpenAI},
  pdfsubject={Multiplicity under arbitrary flat local homomorphisms}}
\setlength{\emergencystretch}{2em}
\setlist{itemsep=3pt,topsep=5pt}
\numberwithin{equation}{section}
\newtheorem{theorem}{Theorem}[section]
\newtheorem{lemma}[theorem]{Lemma}
\newtheorem{proposition}[theorem]{Proposition}
\newtheorem{corollary}[theorem]{Corollary}
\theoremstyle{definition}
\newtheorem{definition}[theorem]{Definition}
\newtheorem{remark}[theorem]{Remark}
\newtheorem{convention}[theorem]{Convention}
\DeclareMathOperator{\length}{\ell}
\DeclareMathOperator{\rank}{rank}
\DeclareMathOperator{\Spec}{Spec}
\DeclareMathOperator{\Proj}{Proj}
\DeclareMathOperator{\gr}{gr}
\DeclareMathOperator{\Tor}{Tor}
\DeclareMathOperator{\Hom}{Hom}
\DeclareMathOperator{\Supp}{Supp}
\DeclareMathOperator{\Frac}{Frac}
\DeclareMathOperator{\pd}{pd}
\DeclareMathOperator{\grade}{grade}
\DeclareMathOperator{\height}{ht}
\DeclareMathOperator{\im}{im}
\DeclareMathOperator{\coker}{coker}
\DeclareMathOperator{\Ann}{Ann}
\DeclareMathOperator{\Cone}{Cone}
\DeclareMathOperator{\Tr}{Tr}
\newcommand{\cO}{\mathcal O}
\newcommand{\PP}{\mathbb P}
\newcommand{\ZZ}{\mathbb Z}
\newcommand{\NN}{\mathbb N}
\newcommand{\QQ}{\mathbb Q}
\newcommand{\RR}{\mathbb R}
\newcommand{\m}{\mathfrak m}
\newcommand{\n}{\mathfrak n}
\newcommand{\p}{\mathfrak p}
\newcommand{\q}{\mathfrak q}
\newcommand{\eHK}{e_{\mathrm{HK}}}
\newcommand{\RGamma}{\mathbf R\Gamma}
\newcommand{\perfd}{\mathrm{perfd}}
\title{Lech's multiplicity conjecture}
\author{OpenAI}
\date{September 23, 2026}

\begin{document}
\maketitle
\begin{abstract}
We prove Lech's multiplicity conjecture: Hilbert--Samuel multiplicity
cannot decrease under a flat local homomorphism of nonzero Noetherian
local rings. The result holds in arbitrary dimension, with no restrictions
on the residue fields or characteristics.
\end{abstract}

\section{Introduction}

For a nonzero Noetherian local ring $(A,\mathfrak a)$ of dimension $d$,
its Hilbert--Samuel multiplicity is
\[
 e(A)=\lim_{N\longrightarrow\infty}
       \frac{d!\,\length_A(A/\mathfrak a^N)}{N^d}.
\]
In dimension zero this is simply the length of $A$.
Multiplicity measures the leading-order growth of the infinitesimal
neighborhoods of the closed point. Lech asked whether this invariant
can decrease under a flat local homomorphism~\cite{Lech,Lech64}.
We prove that it cannot.

\begin{theorem}[Lech's conjecture]\label{thm:main}
Let $(R,\mathfrak m)\longrightarrow(S,\mathfrak n)$ be a flat local
homomorphism of nonzero Noetherian local rings. Then
\[
                         e(R)\leq e(S).
\]
There is no restriction on the dimensions, the residue fields, or the
characteristics of the rings.
\end{theorem}

Flatness compares the quotients by $\mathfrak m^N S$, whereas $e(S)$
is defined by the powers of $\mathfrak n$. Even when the closed fiber
is zero-dimensional, so that both filtrations have finite colength,
this difference prevents a direct comparison of their leading terms
from giving the asserted lower bound for $e(S)$.

\paragraph{Previous work.}
Lech's original work established the inequality when the source has
dimension at most two and in important cases with complete-intersection
closed fiber~\cite{Lech,Lech64}. His original equal-dimensional
formulation leads to the general one by the reductions recalled in
Section~\ref{sec:reductions}.
Ma proved the conjecture for equal-characteristic sources of
dimension three and established the general bound
\[
 e(R)\leq \max\{1,d!/2^d\}\,e(S),\qquad d=\dim R,
\]
in equal characteristic~\cite[Theorems~1.2 and~1.3]{Ma}.
His use of Cohen factorizations and Hilbert--Kunz multiplicity is
also central to the characteristic-$p$ argument here.

\paragraph{Ulrich methods and their limits.}
An Ulrich module is a nonzero finitely generated maximal
Cohen--Macaulay module whose multiplicity equals its minimal number
of generators. Such modules turn multiplicity into a generator count
that can be compared across a flat map. Hanes developed this approach,
including its asymptotic form: maximal Cohen--Macaulay modules of
positive rank whose generator-to-multiplicity ratios tend to one
also suffice. He obtained graded special cases by this method
\cite[Chapter~IV, especially Proposition~4.1.5]{Hanes}.
Ma's weakly lim Ulrich sequences allow the modules themselves to
fail to be Cohen--Macaulay: their multiplicity-to-generator ratios
tend to one, while their first higher parameter Koszul Euler
characteristics are negligible relative to the generator counts
\cite[Definitions~2.2 and~2.9]{MaUlrich}.
For a local domain, existence of such a sequence implies the
multiplicity inequality for every flat local target
\cite[Theorem~2.11]{MaUlrich}.

In arbitrary dimension, Ma proved the conjecture when the source
is a standard graded algebra over a perfect field, localized at
its homogeneous maximal ideal; the flat local target need not be
graded~\cite[Theorem~3.8]{MaUlrich}. The proof constructs weakly lim
Ulrich sequences in positive characteristic and uses specialization
in characteristic zero. Passing to associated graded rings does
not extend this result to every flat map, because the induced map
of associated graded rings need not remain flat
\cite[Remark~3.9]{MaUlrich}.
A complementary graded case was proved by Meng: for an
equal-dimensional flat local map, the inequality holds when, up to
completion, the target is standard graded over a field and the
extended source maximal ideal is homogeneous
\cite[Theorem~5.13]{Meng}.

Nor can one require weakly lim Ulrich sequences over every domain:
Yhee constructed complete two-dimensional local domains admitting
no such sequence~\cite[Theorem~6.2]{Yhee}. These examples are
nonnormal. A different obstruction, within normal domains, is supplied
by a companion construction of a complete three-dimensional Noetherian
normal local $\mathbb C$-domain, with residue field $\mathbb C$, admitting
no nonzero finitely generated maximal Cohen--Macaulay module
\cite[Theorem~1.1]{SmallCM}. Thus approaches requiring such a module
over every source ring itself cannot be universal; this finite-module
obstruction does not rule out weakly lim Ulrich sequences.

\paragraph{From modules to complexes.}
Over a $d$-dimensional local ring, a short complex is a finite free
complex in degrees $0,\ldots,d$ with finite-length homology and
nonzero zeroth homology.
Iyengar, Ma, and Walker developed multiplicity and Betti-number
inequalities for short free complexes using lim Ulrich methods
\cite{IMW2022}. Their approach to Lech's conjecture in
\cite[Proposition~8.3]{IMW2022} uses a Cohen factorization and the
tensor product of a finite free resolution with a parameter Koszul
complex. This is the predecessor of the construction used here.
They also used lim Ulrich sequences to compare cones of Betti tables
with those over linear Noether normalizations~\cite{IMW2023}.
More recently, they related Ulrich modules to extensions of sheaves
from the exceptional fiber of a blowup
\cite[Theorem~1.1]{IMW2025}, obtaining in dimension two
short-complex inequalities under geometric hypotheses
\cite[Theorems~5.2 and~6.9]{IMW2025}.
These results motivate seeking positivity directly for complexes,
without requiring Ulrich modules or sequences over every source.

Two difficulties remain in passing from such results to the
unrestricted theorem: retaining coefficient $1$ in every dimension,
and obtaining the required length estimates in mixed characteristic.
Our argument uses a single estimate for free complexes in both
settings. Its constants depend only on the number of parameters and the ranks,
even when the auxiliary ring or its length function changes.

\paragraph{A uniform estimate for complexes.}
The main technical ingredient is Theorem~\ref{thm:complex}. Let
$I=(z_1,\ldots,z_h)C$ in a possibly non-Noetherian ring $C$, and suppose an
additive length $\lambda$, satisfying the stability hypotheses stated
there, makes the powers of $z$ behave as a parameter
system: their positive Koszul homology has length zero and their quotient
lengths scale as $a^h$. A finite free complex $F$ of length at most $h$
and alternating rank zero, contractible off $V(I)$, whose matrices have
entries in $I^s$, satisfies
\[
 \lambda(H_0F)\ \geq\ \lambda(C/I)s^h
       \left(\rank F_0-\frac{A_h\sum_i\rank F_i}{s}\right)
\]
for sufficiently large $s$. The constant does not depend on the ring,
the length theory, or any auxiliary finite extension.

Here $H_0F$ denotes the zeroth homology of $F$. The hypotheses on
$\lambda$ are part of the estimate: additivity alone does not suffice.
In particular, sums and filtered colimits of modules of length zero
must still have length zero when one power of $I$ annihilates all
the modules. Positive Koszul homology is required to have length
zero, and need not itself vanish.

The proof combines a blowup comparison with the characteristic-free
vector-bundle construction of Eisenbud--Schreyer~\cite{EisenbudSchreyer},
which gives cohomology in a single degree. The relationship between
free complexes, cohomology tables, and positivity has also been
developed in the categorified duality of
Eisenbud--Erman~\cite{EisenbudErman}. Here the needed comparison of Euler
characteristics is proved using finite filtrations and stabilized images
of truncation maps. This comparison is needed because the ring $C$
need not be Noetherian. On the blowup, the differential entries can
be divided by appropriate powers of the invertible ideal $I\cO$.
Tensoring the resulting complex with a suitable vector bundle places
its global hyperhomology in degree zero, modulo length zero, and
hence gives a nonnegative Euler characteristic. The bundle's
cohomological roots are the integer twists at which all its
cohomology vanishes. These roots, after division by $s$, are chosen
near the roots of the derivative
of $\prod_{i=1}^h(x+i)$; this choice makes all
the relevant layer sums have the same leading term. Here the layers
are successive quotients for powers of the invertible ideal
$I\cO$ on the blowup. The alternating
rank identity then leaves precisely the contribution of $F_0$.

The estimate applies to length theories beyond ordinary module length.
For the present problem we use perfection in characteristic $p$ and the
perfectoid normalized lengths of Cai--Lee--Ma--Schwede--Tucker in mixed
characteristic~\cite{CLMST}.

\paragraph{A Dutta-multiplicity consequence.}
In characteristic $p$, the same estimate also applies to arbitrary
short complexes. For such a complex over a $d$-dimensional ring of
characteristic $p>0$, its Dutta multiplicity~\cite{Dutta} is the
Frobenius-normalized Euler
characteristic
\[
 \chi_\infty^R(F)=\lim_{n\to\infty}p^{-nd}
       \sum_i(-1)^i\length_R H_i(\Phi_R^nF),
\]
where $\Phi_R^nF$ is obtained by raising the differential entries to
their $p^n$th powers.
Corollary~\ref{cor:dutta-domain} proves
$\chi_\infty(F)\geq e(R)$ for every short complex over a complete
Noetherian local domain of characteristic $p>0$.  This is the case of
complete local domains of positive characteristic in the conjecture of
Iyengar--Ma--Walker~\cite[Conjecture~8.1]{IMW2022}.
Their full conjecture, for all complete local rings, would imply
Lech's conjecture~\cite[Proposition~8.3]{IMW2022}. Here the stated
domain case is a consequence of the uniform estimate, not an
assumption in the proof of Lech's conjecture.

\paragraph{How the estimate proves the conjecture.}
After reductions that preserve multiplicity or use its additivity,
we work with a complete domain source, equal source and target
dimensions, and an algebraically closed target residue field.
In positive residue characteristic, a Cohen factorization supplies
a complete local ring $T$ and a perfect ideal $J$ such that
$S=T/J$ and $e(T)=e(R)$. If $d=\dim S$ and $h=\dim T$, a finite
free resolution $P$ of $S$ over $T$ has length $h-d$.
Choose parameters $x_1,\ldots,x_d$ in $S$ whose ideal is a reduction
of its maximal ideal, so its multiplicity is $e(S)$, and lift them
to $T$. The complex
\[
 F=P\otimes_T K(x_1,\ldots,x_d;T)
\]
has length $h$, has degree-zero rank one, and is contractible away
from the closed point. Here $K(x;T)$ is the Koszul complex.
This is the complex to which the uniform estimate is applied.

The two positive-residue-characteristic arguments turn this complex
into a multiplicity comparison in different ways. In characteristic
$p$, Frobenius raises the differential entries to arbitrarily high
order without changing their ranks. On each irreducible component
of $T$, normalized length on the perfection identifies the length of
the resulting zeroth homology with a Hilbert--Kunz multiplicity. The complex
estimate gives a lower bound by the Hilbert--Samuel multiplicity of
that component. Associativity and Ma's perfect-ideal identity then
identify the sum with $e(S)$.

In mixed characteristic, roots of parameters and Heitmann's
integral-extension Brian\c{c}on--Skoda theorem~\cite{Heitmann}, in
the form recalled in~\cite[Proposition~2.8]{HeitmannMa}, replace
Frobenius. The extension has a rank that may grow as the
entries become deeper. Dividing normalized lengths by that rank
cancels it from the lower bound. For each fixed extension we obtain
an upper bound by replacing $x_i$ with $x_i^N$ and taking $N$ to
infinity in a prime filtration of $S$. Only after this limit do we
let the depth of the differential entries tend to infinity.
The upper bound is independent of the chosen extension. Moreover,
each top-dimensional filtration prime lies on exactly one component
of $T$, which is the reason this summation retains coefficient one.

In residue characteristic zero, we first replace the given map by
a finite free map with a common coefficient field. A finite set of
polynomial identities records algebra structures, freeness,
locality, a lower bound on the source multiplicity, and an upper
bound on the target multiplicity. The last two conditions use
presentation minors and a reduction relation on the source, and
monic initial relations for the target's associated graded ring.
Artin approximation~\cite{Artin} and specialization preserve these identities
in positive characteristic. A strict counterexample would therefore
contradict the characteristic-$p$ case.

\paragraph{Organization.}
Section~\ref{sec:complex} proves the uniform estimate.
Section~\ref{sec:lengths} constructs its length inputs in positive
and mixed characteristic. Section~\ref{sec:reductions} makes the
flat-local reductions and constructs $T$ and $F$.
Sections~\ref{sec:charp} and~\ref{sec:mixed} prove the two cases of
positive residue characteristic; Section~\ref{sec:charp} also proves
the Dutta-multiplicity corollary. Section~\ref{sec:charzero} gives
the finite encoding and specialization in residue characteristic
zero, and Section~\ref{sec:assembly} assembles the theorem.

\paragraph{Conventions and background.}
Rings are commutative with identity. Local rings are Noetherian
unless explicitly introduced as auxiliary rings for an additive
length; those auxiliary rings remain exempt from this convention
when used in later sections. For an ideal of definition $H$ of a local ring $A$, write
$e(H,A)$ for its Hilbert--Samuel multiplicity. A system of parameters
has length $\dim A$. A \emph{reduction} of an ideal $I$ is a subideal
$Q$ such that $I^{n+1}=QI^n$ for all sufficiently large $n$; it is
\emph{minimal} if it is inclusion-minimal among reductions of $I$.
For an ideal of definition, $e(Q,A)=e(I,A)$. Indeed, the reduction
relation gives $I^{n+c}\subseteq Q^n\subseteq I^n$ for one fixed
$c$ and all $n$, so their leading length coefficients agree.
We use homological indexing: a sheaf-cohomology
degree $v$ contributes total homological degree $i-v$ to a term of
complex degree $i$. The shift $F[t]$ has $H_i(F[t])=H_{i-t}(F)$.
We use the standard additivity and associativity formulas for
multiplicity, and state the deeper literature inputs where they enter.
In particular, the perfectoid length theorems are inputs, not claims
that the auxiliary algebras are flat or Noetherian.

\section{A complex estimate with an additive length}\label{sec:complex}

The goal of this section is to bound the zeroth homology of a short
free complex from below using only the order of its differential
entries and the ranks of its terms. We first isolate the length
hypotheses, then prove the comparison between the complex and its
pullback to a blowup. A vector bundle with prescribed cohomological
roots will turn that comparison into the required inequality.

\subsection{Length hypotheses and the estimate}

All complexes in this section are indexed homologically. Thus a sheaf
cohomology group in degree $v$, arising from the term of a complex in degree
$i$, contributes to total degree $i-v$.

Let $C$ be a commutative ring, let $h\geq 1$, and put
$I=(z_1,\ldots,z_h)C$. Write $\mathscr T_I$ for the abelian category of
$C$-modules annihilated by some power of $I$. Suppose that
\[
 \lambda:\mathscr T_I\longrightarrow [0,\infty]
\]
is additive on short exact sequences and satisfies $\lambda(0)=0$.
We assume, in addition, that direct sums and filtered colimits of
length-zero modules annihilated by a common power of $I$ have length zero.
The standing numerical and homological assumptions are
\begin{equation}\label{eq:complex-length-hypotheses}
 \begin{split}
 0<\mu&:=\lambda(C/I)<\infty,\\
 \lambda\bigl(C/(z_1^a,\ldots,z_h^a)C\bigr)&=a^h\mu
       \qquad(a\geq 1),\\
 \lambda H_i(z_1^a,\ldots,z_h^a;C)&=0
       \qquad(a\geq 1,\ i>0).
 \end{split}
\end{equation}
No Noetherian hypothesis is imposed on $C$.

We say that a map between objects of $\mathscr T_I$ is an
\emph{isomorphism modulo length zero} if its kernel and cokernel have
length zero. The length-zero objects form a Serre subcategory. Consequently
one may work in the quotient by this subcategory when taking kernels,
images, cokernels, or finite spectral sequences. In particular, a
commutative square whose two vertical maps are isomorphisms modulo length
zero induces such an isomorphism between the images of its horizontal
maps. All uses of this convention below involve finite filtrations.

We will also use the following elementary observation about complexes
with finite exhaustive filtrations in a common index range. If a
filtration-preserving map induces zero on the first page of its spectral
sequence, it induces zero on the associated graded of homology and hence
lowers the filtration on homology by one. A composition of as many such
maps as there are filtration steps is therefore zero on homology. The same
statement holds modulo length zero whenever the pages under consideration
belong to $\mathscr T_I$.

\begin{theorem}\label{thm:complex}
Let $C$, $I$, and $\lambda$ be as above. Assume that $\lambda$ is additive
on short exact sequences, that $\lambda(0)=0$, and that direct sums and
filtered colimits of length-zero modules killed by a common power of $I$
have length zero. Assume also \eqref{eq:complex-length-hypotheses}.
Let $F$ be a complex of finite
free $C$-modules, concentrated in degrees $0,\ldots,h$, which is
contractible locally on $\Spec C\setminus V(I)$. Set
$b_i=\rank_C F_i$. Then $H_i(F)$ is annihilated by a power of $I$,
\[
 \lambda H_i(F)=0\quad(i>0),\qquad \lambda H_0(F)<\infty.
\]
Suppose further that
\[
 \sum_{i=0}^h(-1)^i b_i=0
 \quad\text{and that all differential entries belong to }I^s,
 \qquad s\geq1.
\]
There is a constant $A_h$, depending only on $h$, such that
\begin{equation}\label{eq:complex-estimate}
 \lambda H_0(F)\ \geq\
 \mu s^h\left(b_0-\frac{A_h}{s}\sum_{i=0}^h b_i\right).
\end{equation}
In particular, for fixed $h$ and fixed ranks the error is $O_h(1/s)$,
uniformly in $C$, $I$, $\lambda$, and $F$.
\end{theorem}

\subsection{Koszul acyclicity modulo length zero}

We first prove the homology assertion, which does not require the
conditions on ranks or differential entries.

\begin{lemma}\label{lem:acyclicity}
Under the standing length hypotheses, let $F$ be a finite free
complex in degrees $0,\ldots,h$, contractible locally off $V(I)$.
Each $H_i(F)$ is annihilated by
a fixed power of $I$, all its lengths are finite, and its positive
homology has length zero.
\end{lemma}
\begin{proof}
For each $j$, the complex $F[1/z_j]$ is exact. A bounded exact complex of
projective modules is contractible, so clearing the finitely many
denominators in a contracting homotopy shows that some power of $z_j$
acts nullhomotopically on $F$. Choose a common exponent $a_0$.
It follows that $(z_1^{a_0},\ldots,z_h^{a_0})$ annihilates every $H_i(F)$,
and therefore so does $I^{h(a_0-1)+1}$.

Put $K_a=K(z_1^a,\ldots,z_h^a;C)$. If $a\geq a_0$, each multiplication
by $z_j^a$ is nullhomotopic on $F$ and also on its tensor product with any
of the other Koszul factors. Splitting the resulting mapping cones gives
\begin{equation}\label{eq:complex-koszul-splitting}
 F\otimes_C K_a\simeq
 \bigoplus_{T\subseteq\{1,\ldots,h\}}F[|T|].
\end{equation}
The splittings can be chosen to preserve the inclusion of the copy of
$F$ in Koszul degree zero.

The homology of $F\otimes K_a$ is annihilated by
$(z_1^a,\ldots,z_h^a)$, since these multiplications are nullhomotopic on
$K_a$. In the bounded double-complex spectral sequence obtained by taking
Koszul homology first, all positive Koszul rows have length zero by
\eqref{eq:complex-length-hypotheses}. The remaining row has terms
\[
 F_i\otimes_C C/(z_1^a,\ldots,z_h^a),\qquad 0\leq i\leq h,
\]
of finite length. Hence every $H_n(F\otimes K_a)$ has finite length, and
its length is zero when $n>h$. The summand $F[h]$ in
\eqref{eq:complex-koszul-splitting} proves the assertions: $H_0(F)$ is a
summand of $H_h(F\otimes K_a)$, and $H_i(F)$ for $i>0$ is a summand of
$H_{h+i}(F\otimes K_a)$.
\end{proof}

\subsection{The blowup and its exceptional divisor}

The homology lengths are now finite. To estimate them, we pass to a
space on which $I$ becomes invertible, while keeping track of the
Euler characteristic of the original complex. Let
\[
 Y=\Proj_C\Bigl(\bigoplus_{n\geq0}I^n\Bigr),\qquad
 L=I\cO_Y=\cO_Y(1),\qquad E=V(L).
\]
Here $\cO_Y(1)$ is the Rees-algebra twist. Indeed, the standard chart
corresponding to $z_i$ has ring $C[I/z_i]\subseteq C[1/z_i]$. On this
chart $I$ becomes $(z_i)$ and $z_i$ is a nonzerodivisor. Thus $L$ is an
invertible ideal, and $E$ is an effective Cartier divisor, even when $C$
is not Noetherian. Empty charts may simply be omitted. The generators of
$I$ give a closed immersion $Y\longrightarrow\PP_C^{h-1}$, and
\[
 E=\Proj(\gr_I C)\longrightarrow\PP_{C/I}^{h-1}.
\]
For a vector bundle $V$ on $\PP_\ZZ^{h-1}$ we use the same symbol for its
pullbacks to these schemes, and write $V(t)=V\otimes\cO(t)$.

\begin{lemma}\label{lem:complex-graded-comparison}
The graded surjection
\begin{equation}\label{eq:complex-graded-map}
 (C/I)[Z_1,\ldots,Z_h]\longrightarrow\gr_I C
\end{equation}
has length-zero kernel in every degree. For every vector bundle $V$ on
$\PP_\ZZ^{h-1}$ and every integer $t$, restriction induces isomorphisms
modulo length zero
\[
 H^v(\PP_{C/I}^{h-1},V(t))\longrightarrow H^v(E,V(t)).
\]
All of these cohomology modules have finite length. Moreover, for $t\geq0$
the natural map
\[
 \gr_I^t C\longrightarrow H^0(E,\cO_E(t))
\]
is an isomorphism modulo length zero, and $H^v(E,\cO_E(t))$ has length
zero for $v>0$.
\end{lemma}
\begin{proof}
Set $J_a=(z_1^a,\ldots,z_h^a)$. Since
$I^{h(a-1)+1}\subseteq J_a$, the $I$-adic filtration of $C/J_a$ is finite.
There is a graded surjection
\[
 (C/I)[Z_1,\ldots,Z_h]/(Z_1^a,\ldots,Z_h^a)
 \longrightarrow\gr_I(C/J_a).
\]
The sum of the lengths of the homogeneous pieces on each side is $a^h\mu$.
Additivity and nonnegativity imply that every homogeneous piece of its
kernel has length zero. For a given degree $n$, choose $a>n$. The relations
$Z_i^a$ have no effect in degree $n$, and $J_a\subseteq I^{n+1}$, so the
degree-$n$ map is precisely \eqref{eq:complex-graded-map}. This proves the
first assertion. It also gives
\begin{equation}\label{eq:complex-power-colength}
 \lambda(C/I^M)=\mu\binom{M+h-1}{h}\qquad(M\geq1).
\end{equation}

Write $P=(C/I)[Z_1,\ldots,Z_h]$ and let $N$ be the kernel in
\eqref{eq:complex-graded-map}. The degree-zero component of a homogeneous
localization of $N$, as used on any standard projective chart or its
intersections, is a filtered colimit of homogeneous pieces of $N$.
All these pieces are annihilated by $I$ and have length zero. The assumed
colimit property therefore gives length zero for every such localized
component. Tensoring with the restriction of $V(t)$ preserves this
property, since that restriction is finite projective over the chart ring.
The finite standard affine Cech complexes for projective space and $E$
consequently have termwise length-zero kernels, and restriction is an
isomorphism modulo length zero on their cohomology.

To check finiteness without a hypothesis on $C/I$, resolve $V(t)$ on
$\PP_\ZZ^{h-1}$ by finitely many finite sums of twists of the structure
sheaf. Such a resolution exists by taking a finite graded free resolution
of a finitely generated graded module representing $V(t)$ over
$\ZZ[Z_1,\ldots,Z_h]$. This polynomial ring is regular of finite
dimension and therefore has finite global dimension
\cite[Tags~07NF and~00OE]{Stacks};
its finitely generated graded projectives are graded free, as follows by
lifting a homogeneous basis modulo $(Z_1,\ldots,Z_h)$ and applying graded
Nakayama. The sheaf resolution is locally split because its final cokernel
is a vector bundle, so it remains exact after arbitrary base change.
The cohomology of every twist on $\PP_{C/I}^{h-1}$ is finite free over
$C/I$. A finite spectral sequence for the resolution proves the required
finiteness. Finally, for $t\geq0$ the polynomial degree-$t$ piece is the
space of global sections of $\cO(t)$ and its higher cohomology vanishes.
Combining these facts with \eqref{eq:complex-graded-map} proves the last
assertions.
\end{proof}

\subsection{Euler characteristics on the blowup}

We now compare the homology lengths of $F$ with those of its pullback
to $Y$. We will also show that, when the alternating rank is zero,
tensoring the pullback with a vector bundle multiplies its Euler
characteristic by the bundle's rank.

For a bounded complex $D$ of quasi-coherent sheaves on $Y$, write
\[
 \mathbb H_k(Y,D)=H_k\mathbf R\Gamma(Y,D).
\]
These groups can be calculated using the finite standard affine Cech
cover, whose intersections are affine. Whenever they are $I$-power torsion
of finite length, put
\[
 \chi_\lambda(D)=\sum_k(-1)^k\lambda\mathbb H_k(Y,D).
\]
Only finitely many degrees occur in our applications.

Pull $F$ back to $Y$ and denote the resulting complex by $F_Y$. If the
differential entries belong to $I^s$, define a complex $G$ by
\begin{equation}\label{eq:complex-G}
 G_i=L^{-is}\otimes_C F_i=\cO_Y(-is)^{b_i}.
\end{equation}
Its differential is induced by that of $F_Y$ after inverting $L$.
On a chart where $L=(z_j)$, the matrix in the bases
$z_j^{-is}F_i$ is $z_j^{-s}d_i$, where $d_i$ is the matrix for $F$.
It is regular because the entries of $d_i$ belong to $L^s$.
The natural inclusions $F_i\to L^{-is}\otimes_C F_i$ give an
inclusion of complexes $F_Y\subseteq G$.

In degree $i$, the quotient $G/F_Y$ has successive layers
$\cO_E(l)^{b_i}$ for $l=-is,\ldots,-1$. When the alternating rank is
zero, these finite layer ranges will account, after the Euler comparison,
for the entire difference between $\chi_\lambda(G\otimes V)$ and
$(\rank V)\lambda H_0(F)$. We first justify the comparison for an
arbitrary bundle $V$; the choice of $V$ that gives a nonnegative Euler
characteristic will follow.

\begin{lemma}\label{lem:complex-twist-maps}
Let $D$ be either $F_Y\otimes V$ or $G\otimes V$, with $V$ a vector
bundle pulled back from $\PP_\ZZ^{h-1}$. There is an integer $M_D$ such
that, for every integer $j$ and every $M\geq M_D$, the inclusion
\[
 D(j+M)\longrightarrow D(j)
\]
induces zero on hyperhomology. All hyperhomology modules of $D(j)$ are
annihilated by a power of $I$ and have finite length.
\end{lemma}
\begin{proof}
On a standard affine chart $U_i$, $L$ is generated by the nonzerodivisor
$z_i$. The chart complex of $D$ consists of finite projective modules and
becomes contractible after inverting $z_i$: away from $E$, both $F_Y$ and
$G$ identify with the pullback of $F$. Clearing denominators in a
contracting homotopy gives an exponent $m_i$ for which multiplication by
$z_i^{m_i}$ is nullhomotopic on this chart. Trivializing $L$ shows that the
same exponent works for every twist. It also works on all intersections
contained in this chart. Choose $m$ at least as large as every $m_i$.

The map $D(j+m)\to D(j)$ is zero on the internal-homology page of the
Cech double complex, since on each intersection it is identified with a
nullhomotopic multiplication by a local equation to the power $m$.
The Cech filtration has at most $h$ steps. The filtered-complex observation
above therefore shows that $D(j+hm)\to D(j)$ is zero on hyperhomology.
Every larger twist map factors through one of these, so $M_D=hm$ works.

For $a\in I^{M_D}$, multiplication by $a$ on $D(j)$ factors as
\[
 D(j)\xrightarrow{a}D(j+M_D)\longrightarrow D(j).
\]
Thus $I^{M_D}$ annihilates its hyperhomology. If $M\geq M_D$, the long
exact sequence for
\[
 0\longrightarrow D(j+M)\longrightarrow D(j)
 \longrightarrow D(j)/D(j+M)\longrightarrow0
\]
injects $\mathbb H_k(Y,D(j))$ into the hyperhomology of the quotient.
That quotient has a finite filtration by $E$-layers, each of whose terms
is a finite sum of bundles $V(t)|_E$. Their sheaf cohomology has finite
length by Lemma~\ref{lem:complex-graded-comparison}. The two finite
spectral sequences, first for each layer and then for the filtration,
give finite length for the quotient's hyperhomology and hence for that
of $D(j)$.
\end{proof}

The Euler characteristics of finite thickenings need not recover that of $F$.
For example, in the alternating-rank-zero case needed for the estimate,
\[
 \sum_k(-1)^k\lambda H_k(F/I^M F)
 =\lambda(C/I^M)\sum_i(-1)^i b_i=0.
\]
Instead, we compare the images of homology maps from a sufficiently
large thickening to a fixed smaller one. These images will recover the
homology of $F$ and of $F_Y$, modulo length zero. This argument gives
the first assertion below without the alternating-rank hypothesis;
that hypothesis is used only afterward to prove invariance under twists
and the vector-bundle formula.

\begin{lemma}\label{lem:complex-euler-comparison}
One has
\begin{equation}\label{eq:complex-euler-comparison}
 \chi_\lambda(F_Y)=\sum_i(-1)^i\lambda H_i(F)=\lambda H_0(F).
\end{equation}
If $\sum_i(-1)^i b_i=0$, then for every vector bundle $V$ pulled back from
$\PP_\ZZ^{h-1}$,
\begin{equation}\label{eq:complex-bundle-rank}
 \chi_\lambda(F_Y\otimes V)=(\rank V)\chi_\lambda(F_Y).
\end{equation}
\end{lemma}
\begin{proof}
\emph{Comparison on finite thickenings.}
For each $M\geq1$ there is a natural comparison
\begin{equation}\label{eq:complex-thick-comparison}
 F/I^M F\longrightarrow
 \mathbf R\Gamma(Y,F_Y/L^M F_Y).
\end{equation}
Filter its two sides by powers of $I$ and $L$, respectively. At the
$l$th layer, $0\leq l<M$, comparison for each free summand is
\[
 \gr_I^l C\longrightarrow\mathbf R\Gamma(E,\cO_E(l)).
\]
Lemma~\ref{lem:complex-graded-comparison} says that this is an isomorphism
on homology modulo length zero. The finite filtration, together with the
finite complex $F$, proves the same assertion for
\eqref{eq:complex-thick-comparison}. These comparisons commute with all
transition maps as $M$ increases.

\smallskip
\noindent\emph{The geometric tower.}
We identify the eventual images of these two towers. On the geometric
side put
\[
 A=\mathbf R\Gamma(Y,F_Y),\qquad
 B_M=\mathbf R\Gamma(Y,F_Y/L^M F_Y).
\]
The triangles
\begin{equation}\label{eq:complex-thick-triangle}
 A\longrightarrow B_M\longrightarrow
 \mathbf R\Gamma(Y,F_Y\otimes L^M)[1]\longrightarrow A[1]
\end{equation}
are compatible with the transition $N\to M$ for $N\geq M$.
Lemma~\ref{lem:complex-twist-maps} implies that $H_k(A)\to H_k(B_M)$ is
injective once $M$ is sufficiently large. It also says that the transition
on the third terms of \eqref{eq:complex-thick-triangle} is zero on homology
once $N-M$ is sufficiently large. Exactness then gives
\begin{equation}\label{eq:complex-geometric-image}
 \operatorname{im}\bigl(H_k(B_N)\to H_k(B_M)\bigr)
   =\operatorname{im}\bigl(H_k(A)\to H_k(B_M)\bigr)
   \cong H_k(A).
\end{equation}
The reverse inclusion in this equality uses only the compatible maps from
$A$ to all $B_N$.

\smallskip
\noindent\emph{The Koszul tower.}
For the algebraic side first use $J_a=(z_1^a,\ldots,z_h^a)$ and
$A_a=F\otimes K_a$. The Koszul augmentation gives a compatible map
\[
 A_a\longrightarrow F/J_aF
\]
which is an isomorphism on homology modulo length zero. Indeed its cone
is obtained by tensoring the bounded free complex $F$ with the cone of
$K_a\to C/J_a$; all homology of the latter cone has length zero by
\eqref{eq:complex-length-hypotheses}.

The inclusion $F\subseteq A_a$ is the subcomplex in Koszul degree zero.
For $a\geq a_0$ it is split on homology by
\eqref{eq:complex-koszul-splitting}. Let $Q_a=A_a/F$. The standard
transition $K_b\to K_a$, for $b\geq a$, is the identity in degree zero
and multiplies the basis vector indexed by $T$ by
$\prod_{j\in T}z_j^{b-a}$. Filter $Q_a$ increasingly by Koszul
degrees $1,\ldots,p$, for $p=1,\ldots,h$.
Its associated graded complexes are copies of shifts of $F$, and the
transition on each such piece has a nonempty index set $T$.
Choose $t$ with $I^tH(F)=0$, as in Lemma~\ref{lem:acyclicity}.
For $b-a\geq t$ every transition is zero on the homology of these graded
pieces. A composition of $h$ such transitions is therefore zero on
$H(Q_a)$. For example, $H(Q_b)\to H(Q_a)$ is zero when $b-a\geq ht$.
The long exact sequences for $0\to F\to A_a\to Q_a\to0$ now give,
for $a$ sufficiently large and then $b$ sufficiently large,
\[
 \operatorname{im}\bigl(H_k(A_b)\to H_k(A_a)\bigr)
  =\operatorname{im}\bigl(H_k(F)\to H_k(A_a)\bigr)
  \cong H_k(F).
\]
Passing through the Koszul augmentations proves this assertion modulo
length zero for the tower $F/J_aF$.

\smallskip
\noindent\emph{From Koszul powers to ordinary powers.}
Here is explicitly why it passes to the ordinary-power tower
$P_M=F/I^M F$. The containments
\begin{equation}\label{eq:complex-cofinality}
 I^{h(a-1)+1}\subseteq J_a\subseteq I^a
\end{equation}
give all the required quotient maps. Choose one sufficiently large $a$.
For $M\geq h(a-1)+1$, the composite
$H_k(F)\to H_k(P_M)\to H_k(F/J_aF)$ proves injectivity of the first
map modulo length zero. Next fix such an $M$, choose $a\geq M$ large,
and choose $b\geq a$ so that the image from level $b$ to level $a$ in
the $J$-tower is the image of $H_k(F)$. For $N\geq h(b-1)+1$, the
transition factors through
\[
 P_N\longrightarrow F/J_bF\longrightarrow F/J_aF\longrightarrow P_M.
\]
Its homology image is consequently contained in the image of $H_k(F)$
modulo length zero; compatibility gives the opposite containment. Thus
the eventual image in the $P$-tower identifies with $H_k(F)$ modulo
length zero.

Choose $M$ and then $N$ large enough for both this description and
\eqref{eq:complex-geometric-image}. The compatible comparisons
$P_N\to B_N$ and $P_M\to B_M$ in
\eqref{eq:complex-thick-comparison} identify the homology images of
$P_N\to P_M$ and $B_N\to B_M$ modulo length zero. These images
identify with $H_k(F)$ and $H_k(A)$, respectively, modulo length zero.
All their lengths are finite by
Lemma~\ref{lem:acyclicity}, \eqref{eq:complex-power-colength}, and
Lemma~\ref{lem:complex-twist-maps}. Therefore
\[
 \lambda\mathbb H_k(Y,F_Y)=\lambda H_k(F)
\]
for every $k$, which proves \eqref{eq:complex-euler-comparison}.

\smallskip
\noindent\emph{Tensoring by a vector bundle.}
The homology comparison is established. To prove
\eqref{eq:complex-bundle-rank}, first compare consecutive twists
of $F_Y$. The quotient of $F_Y(j)$ by $F_Y(j+1)$ has term
$\cO_E(j)^{b_i}$ in degree $i$, so its Euler characteristic is
\[
 \left(\sum_i(-1)^i b_i\right)
 \sum_v(-1)^v\lambda H^v(E,\cO_E(j))=0.
\]
Euler additivity shows that $\chi_\lambda(F_Y(j))$ is independent of
$j\in\ZZ$. Finally resolve $V$ on $\PP_\ZZ^{h-1}$ by a finite complex
of finite sums of twists, as in
Lemma~\ref{lem:complex-graded-comparison}. The resolution remains exact
on $Y$. Tensoring with $F_Y$ and taking Euler characteristics expresses
$\chi_\lambda(F_Y\otimes V)$ as the alternating sum of the corresponding
twist characteristics. The alternating sum of the ranks of the resolution
is $\rank V$, giving \eqref{eq:complex-bundle-rank}.
\end{proof}

\subsection{Bundles with prescribed cohomological roots}

We next construct the bundle that determines the signs of the needed
Euler characteristic. We use the product-of-projective-lines
construction of Eisenbud--Schreyer
\cite[Theorem~6.1]{EisenbudSchreyer}, verifying its integral and
base-change properties so that it applies over our auxiliary ring.

\begin{lemma}\label{lem:complex-root-bundle}
Given integers $r_1>\cdots>r_{h-1}$, there is a vector bundle $V$ of rank
$(h-1)!$ on $\PP_\ZZ^{h-1}$ with the following property. For every ring
$A$ and every integer $a$, $V(a)$ on $\PP_A^{h-1}$ is acyclic if $a=r_j$
for some $j$. Otherwise its cohomology is a finite free $A$-module in
the single degree
\[
 v(a)=\#\{j:r_j>a\},
\]
and its Euler characteristic in ranks is
\[
 \prod_{j=1}^{h-1}(a-r_j).
\]
For $A=C/I$, the Euler characteristic with respect to $\lambda$ is
$\mu\prod_j(a-r_j)$, and the same Euler characteristic and cohomology
vanishing statements hold on $E$, modulo length zero.
\end{lemma}
\begin{proof}
Put $n=h-1$. If $n=0$, use the rank-one bundle on $\PP_\ZZ^0$.
For $n>0$, multiplication of binary linear forms defines
\[
 f:(\PP_\ZZ^1)^n\longrightarrow\PP_\ZZ^n,
 \qquad f^*\cO(1)=\cO(1,\ldots,1).
\]
The map is proper and has finite geometric fibers, since a binary form
has only finitely many ordered factorizations into linear factors, up to
the scalars already removed by the projective coordinates. Hence $f$ is
finite. On every field fiber, source and target are smooth of dimension
$n$. Corresponding local rings have equal dimensions, since the residue
field extension is finite and the dimensions are $n$ minus the respective
residue-field transcendence degrees. Thus the source is Cohen--Macaulay
and satisfies the dimension criterion for flatness over the regular
target~\cite[Tag~00R4]{Stacks}. This gives flatness on each field fiber.
Since both schemes are
flat over $\ZZ$, the
fiberwise flatness criterion~\cite[Tag~039D]{Stacks} gives flatness
of $f$. Over $\QQ$ the finite
extension of function fields is separable. A geometric generic form has
$n$ distinct factors and exactly $n!$ ordered factorizations. Its reduced
geometric generic fiber therefore has $n!$ points, so the generic degree,
and hence the constant rank of $f$, is $n!$.

Define
\[
 V=f_*\cO(-1-r_1,\ldots,-1-r_n).
\]
Finite flatness makes this a vector bundle of rank $n!$, and its formation
commutes with arbitrary base change. The projection formula reduces the
cohomology of $V(a)$ to that of
$\cO(a-1-r_1,\ldots,a-1-r_n)$ on $(\PP_A^1)^n$.
On $\PP_A^1$, the bundle $\cO(m)$ has free cohomology only in degree zero
for $m\geq0$, only in degree one for $m\leq-2$, and none for $m=-1$;
its Euler characteristic is $m+1$. These computations commute with any
base change and their complexes split over $\ZZ$. Tensoring the complexes
for the $n$ factors proves the asserted cohomological degree and product
formula over every $A$. For $A=C/I$, a free module of rank $q$ has length
$q\mu$. The statements on $E$ follow from
Lemma~\ref{lem:complex-graded-comparison}.
\end{proof}

\subsection{The sign and the leading coefficient}

\begin{proof}[Proof of Theorem~\ref{thm:complex}]
The homology assertions are Lemma~\ref{lem:acyclicity}. It remains to
prove the estimate. We choose the bundle roots to serve two purposes:
their positions force the hyperhomology of $G\otimes V$ into degree zero,
modulo length zero, while their limiting values make every layer sum
with $1\leq i\leq h$ have the same leading term $(h-1)!s^h$.
The alternating-rank identity will then isolate $b_0$.
The derivative of the following polynomial provides exactly these roots:
\[
 H(x)=\prod_{i=1}^h(x+i),
\]
and write $\rho_1>\cdots>\rho_{h-1}$ for the roots of $H'(x)$.
Rolle's Theorem gives
\[
 -(j+1)<\rho_j<-j.
\]
For all sufficiently large $s$, depending only on $h$, round $s\rho_j$
to integers $r_j$ so that
\begin{equation}\label{eq:complex-root-positions}
 -(j+1)s<r_j<-js,\qquad |r_j-s\rho_j|\leq1.
\end{equation}
Let $V$ be the bundle of Lemma~\ref{lem:complex-root-bundle}, and put
$D=G\otimes V$ with $G$ as in \eqref{eq:complex-G}.

We claim that
\begin{equation}\label{eq:complex-positive-chi}
 \chi_\lambda(D)\geq0.
\end{equation}
In the $l$th $E$-layer of a twist of $D$, the term in homological degree
$i$ is $V(l-is)|_E^{b_i}$. If $l\geq0$, then
\eqref{eq:complex-root-positions} implies
\[
 \#\{j:r_j>l-is\}\leq i.
\]
Indeed a root with index $j\geq i$ is strictly less than $-is$, and
for $i=0$ all roots are negative. Thus the layer's sheaf cohomology in
degree $v>i$ has length zero. In the finite hyperhomology spectral
sequence no term of negative total degree $i-v$ survives modulo length
zero. Filtering by the layers $l=0,\ldots,M-1$ shows that
\[
 \lambda\mathbb H_k(Y,D/D(M))=0\qquad(k<0).
\]
For $M$ large, Lemma~\ref{lem:complex-twist-maps} makes the map on
hyperhomology from $D(M)$ to $D$ zero. The quotient long exact sequence
therefore injects $\mathbb H_k(Y,D)$ into
$\mathbb H_k(Y,D/D(M))$, proving length zero for $k<0$.

If $l<0$, then
\[
 \#\{j:r_j>l-is\}\geq i-1.
\]
For $i\geq1$ this follows because every $j\leq i-1$ has
$r_j>-(j+1)s\geq-is>l-is$; for $i=0$ the assertion is vacuous.
Consequently the corresponding layer has hyperhomology of length zero
in total degrees greater than $1$. Filtering the quotient with layers
$l=-M,\ldots,-1$ gives
\[
 \lambda\mathbb H_k(Y,D(-M)/D)=0\qquad(k>1).
\]
For large $M$ the map from $D$ to $D(-M)$ is zero on hyperhomology.
The long exact sequence gives a surjection
\[
 \mathbb H_{k+1}(Y,D(-M)/D)\longrightarrow\mathbb H_k(Y,D),
\]
and hence length zero for $k>0$. Lemma~\ref{lem:complex-twist-maps}
supplies finite length in all degrees. We have proved
$\chi_\lambda(D)=\lambda\mathbb H_0(Y,D)\geq0$, as claimed.

The inclusion $F_Y\otimes V\subseteq G\otimes V$ has, in degree $i$,
quotient $(L^{-is}/\cO_Y)\otimes V^{b_i}$, with successive layers
$V(l)|_E^{b_i}$ for $l=-is,\ldots,-1$. All Euler characteristics here
are finite. Additivity, Lemma~\ref{lem:complex-euler-comparison}, and
Lemma~\ref{lem:complex-root-bundle} give
\begin{equation}\label{eq:complex-sum-estimate}
 0\leq (h-1)!\lambda H_0(F)
 +\mu\sum_{i=0}^h(-1)^i b_i
       \sum_{l=-is}^{-1}\prod_{j=1}^{h-1}(l-r_j).
\end{equation}
The sum for $i=0$ is empty.

We make the error in these sums explicit. Set
\[
 P_s(x)=\prod_{j=1}^{h-1}(x-r_j/s),\qquad
 P(x)=\prod_{j=1}^{h-1}(x-\rho_j)=\frac{H'(x)}h.
\]
On the fixed interval $[-h,0]$, \eqref{eq:complex-root-positions}
implies $\|P_s-P\|_\infty=O_h(s^{-1})$, while $P$ and its derivative
are bounded in terms of $h$ alone. The elementary error estimate for
left-endpoint Riemann sums consequently gives, uniformly for
$i=1,\ldots,h$,
\begin{align*}
 s^{-h}\sum_{l=-is}^{-1}\prod_j(l-r_j)
 &=\frac1s\sum_{l=-is}^{-1}P_s(l/s)\\
 &=\int_{-i}^0P(x)\,dx+O_h(s^{-1})\\
 &=\frac{H(0)-H(-i)}h+O_h(s^{-1})
  =(h-1)!+O_h(s^{-1}).
\end{align*}
Since $\sum_{i=1}^h(-1)^i b_i=-b_0$, substitution into
\eqref{eq:complex-sum-estimate} yields
\[
 \lambda H_0(F)\geq
 \mu s^h\left(b_0-O_h\left(\frac{\sum_i b_i}{s}\right)\right).
\]
The error bound and the threshold for $s$ depend only on $h$.
The homotopy exponents and truncation levels used earlier may depend on
$C$, $F$, and $V$; they establish exact comparisons and vanishing modulo
length zero and do not enter this numerical error. This proves
\eqref{eq:complex-estimate} for all sufficiently large $s$. Increase
$A_h$, if necessary, to be at least the threshold for $s$. For the
remaining positive integers $s$, its right side is nonpositive because
$b_0\leq\sum_i b_i$, so nonnegativity of length proves the same inequality.
For $h=1$ the empty-product conventions give the argument directly,
with exact sums and no error. The proof is complete.
\end{proof}

\section{Normalized length algebras}
\label{sec:lengths}

We establish the length inputs for Theorem~\ref{thm:complex} and
Lemma~\ref{lem:acyclicity}.  The mixed-characteristic input uses the
perfectoid results of Cai--Lee--Ma--Schwede--Tucker.  We give the
characteristic-$p$ argument in detail, including the passage to the
non-Noetherian perfection.

For a Noetherian local ring $D$ of characteristic $p$ and a primary
ideal of definition $H$, write $H^{[q]}$ for the ideal generated by the
$q$th powers of elements of $H$, where $q=p^n$. Its Hilbert--Kunz
multiplicity is
\[
 e_{\mathrm{HK}}(H,D)
   =\lim_{q=p^n\longrightarrow\infty}
        \frac{\length_D(D/H^{[q]})}{q^{\dim D}}.
\]
The existence of this limit is Monsky's theorem~\cite{Monsky}; see
also \cite[Section~2.2]{Ma}.

\begin{theorem}\label{thm:length-inputs}
Let $(D,\mathfrak d,k)$ be a complete Noetherian local domain of dimension
$h>0$, with perfect residue field $k$ of characteristic $p>0$.
There exist a $D$-algebra $C$ and a function $\lambda$, with values in
$[0,\infty]$, on $C$-modules annihilated by some power of $\mathfrak d$,
with the following properties.
\begin{enumerate}
\item The function is additive on short exact sequences and vanishes on
the zero module.  Direct sums, sums, and directed colimits of modules of
length zero have length zero whenever a fixed power of $\mathfrak d$
annihilates all modules under consideration.
\item For every system of parameters $z_1,\ldots,z_h$ of $D$ and every
integer $a\geq 1$,
\[
 \lambda\bigl(C/(z_1^a,\ldots,z_h^a)C\bigr)
   =a^h e((z_1,\ldots,z_h),D),
 \qquad
 \lambda H_i(z_1^a,\ldots,z_h^a;C)=0\quad(i>0).
\]
In particular, the number
$\mu=\lambda(C/(z_1,\ldots,z_h)C)$ is positive and finite.
\item If $H$ is $\mathfrak d$-primary, then
\[
 \begin{cases}
 \lambda(C/HC)=e_{HK}(H,D),&\operatorname{char}D=p,\\
 \lambda(C/HC)\leq e(H,D),&\operatorname{char}D=0.
 \end{cases}
\]
\end{enumerate}
The same zero-length stability applies to modules killed by a fixed
power of an ideal extended from any $\mathfrak d$-primary ideal of $D$.
\end{theorem}

For $I=(z_1,\ldots,z_h)C$, assertions (1) and (2) supply the
stability, parameter-length, and Koszul hypotheses of
Theorem~\ref{thm:complex}, with $\mu=e((z),D)$.
Assertion (3) then connects the resulting zeroth-homology bound to
ordinary multiplicities: it identifies the quotient length with
Hilbert--Kunz multiplicity in characteristic $p$, and bounds it by
Hilbert--Samuel multiplicity in mixed characteristic.

\subsection{The normalized length formalism}

We use Faltings' normalized length~\cite[Appendix~2]{Faltings},
developed further by Shimomoto~\cite[Section~2]{Shimomoto}, in the
formulation of \cite[Section~2.2]{CLMST}. Let $(A,\mathfrak
m_A,k)$ be one of the complete regular local rings
\[
 A=k[[t_1,\ldots,t_h]]
 \quad\hbox{or}\quad
 A=W(k)[[t_2,\ldots,t_h]],
\]
and in the second case put $t_1=p$.  Choose compatible $p$-power roots
and set
\[
 A_n=A[t_1^{1/p^n},\ldots,t_h^{1/p^n}],
 \qquad \Lambda=\varinjlim_n A_n.
\]
Thus $\Lambda=A_{\mathrm{perf}}$ in characteristic $p$.  In mixed
characteristic $\Lambda$ denotes the uncompleted tower; its $p$-adic
completion is denoted by $A_\infty$.

We consider the abelian category of $\Lambda$-modules annihilated by
$\mathfrak m_A^N$ for some integer $N$ depending on the module.  The
annihilation here is uniform over the whole module.  In both constructions
below we choose $A\subseteq D$ with $\mathfrak m_A\subseteq\mathfrak d$
and evaluate lengths by restriction of scalars.  Thus every $C$-module killed
by a power of $\mathfrak d$ belongs to this category.
A finitely presented
module in this category descends to a finite-length $A_n$-module $M_n$,
and its normalized length is
\begin{equation}\label{eq:normalized-finite-stage}
 \lambda_\infty(M_n\otimes_{A_n}\Lambda)
       =p^{-nh}\length_{A_n}(M_n).
\end{equation}
The use of $A_n$-length instead of $A$-length makes no difference,
because the two rings have the same residue field.  For a finitely
generated module the length is the infimum of the lengths of finitely
presented modules surjecting onto it.  For an arbitrary module it is
then defined by
\begin{equation}\label{eq:normalized-supremum}
 \lambda_\infty(M)=
 \sup\{\lambda_\infty(N):N\subseteq M
                      \text{ is finitely generated over }\Lambda\}.
\end{equation}
These definitions are compatible and give an additive function on
short exact sequences in this category
\cite[Definition~2.2.1 and Propositions~2.2.3, 2.2.6(a)]{CLMST}.
In particular, length is monotone under submodules and quotients.

We record explicitly the consequences for infinite modules that we
will need.

\begin{lemma}\label{lem:normalized-limit}
Suppose that all the modules in the following assertions are
annihilated by a fixed power of $\mathfrak m_A$.
\begin{enumerate}
\item Direct sums, sums, and directed colimits of zero-length modules
have length zero.
\item If $M=\bigcup_n M_n$ is an increasing union of submodules, then
$\lambda_\infty(M)=\sup_n\lambda_\infty(M_n)$.
\item If $M=\varinjlim_n M_n$ and
$\lambda_\infty(M_n)\longrightarrow 0$, then
$\lambda_\infty(M)=0$.  No injectivity assumption on the transition maps
is required in this assertion.
\end{enumerate}
\end{lemma}

\begin{proof}
A finitely generated submodule of a direct sum is contained in a
finite subsum.  If the summands have length zero, additivity gives
length zero for each finite subsum, and
\eqref{eq:normalized-supremum} gives length zero for the direct sum.
Both a sum of submodules and a directed colimit are quotients of
direct sums, proving the first assertion.  The common annihilating
power ensures that those direct sums belong to the length category.

For the second assertion, each finitely generated submodule of $M$ is
contained in some $M_n$.  Apply \eqref{eq:normalized-supremum} and
monotonicity.

For the third assertion, let $N\subseteq M$ be finitely generated over
$\Lambda$.  Representatives of a finite generating set lie in a common
stage $M_{n_0}$.  Consequently $N$ is contained in the image of $M_n$
in $M$ for every $n\geq n_0$.  Monotonicity for submodules and quotients
gives
\[
 0\leq\lambda_\infty(N)\leq\lambda_\infty(M_n)
       \qquad(n\geq n_0).
\]
The right side tends to zero.  Now take the supremum over $N$.
\end{proof}

\subsection{Characteristic \texorpdfstring{$p$}{p}}

Assume first that $D$ has characteristic $p$.  The Cohen--Gabber
normalization theorem~\cite[Theorem~1.1]{KuranoShimomoto} provides
a coefficient field $k\subseteq D$ and
a system of parameters $t_1,\ldots,t_h$ such that
\[
 A=k[[t_1,\ldots,t_h]]\subseteq D
\]
is module-finite and generically separable; this is also the setup of
\cite[Section~3]{CLMST}.  Work inside a fixed algebraic closure of
$\operatorname{Frac}(D)$ and put
\[
 A'=A_{\mathrm{perf}},\qquad C=D_{\mathrm{perf}},\qquad
 K=\operatorname{Frac}(A),\quad L=\operatorname{Frac}(D),\quad
 K'=K_{\mathrm{perf}}=\operatorname{Frac}(A').
\]
We apply $\lambda_\infty$ to $C$-modules by restriction of scalars to
$A'$.  To compute these lengths, we approximate $C$ by algebras
obtained from finite root stages of $D$.

For every $q=p^n$, define
\[
 B_q=A'\otimes_{A^{1/q}}D^{1/q}.
\]
The natural map from this tensor product to $C$ is injective.  Indeed,
$D^{1/q}$ is torsion-free over $A^{1/q}$, and $A'$ is flat over
$A^{1/q}$: the intervening extensions of regular power-series rings
are free.  Thus $B_q$ injects into
\[
 K'\otimes_{K^{1/q}}L^{1/q}.
\]
The latter ring is a field, because the finite separable extension
$L^{1/q}/K^{1/q}$ is linearly disjoint from the purely inseparable
extension $K'/K^{1/q}$.  This identifies $B_q$ with the subring
$A'D^{1/q}$ of $C$.  These subrings increase with $q$ and exhaust $C$:
each contains $D^{1/q}$, and every finite collection of their
coefficients lies in a sufficiently high root extension of $D$.

We give a self-contained form of the finite-stage comparison in
\cite[Section~3 and Theorem~3.0.7]{CLMST}, beginning with a trace-dual
proof of the conductor bound from \cite[Lemma~2.3]{PolstraTucker}.
Although $B_q$ injects into $C$, passing to a primary quotient can
create a kernel.  The conductor bound below controls this kernel by
giving elements that multiply $C$ into $B_q$.

\begin{lemma}\label{lem:perfection-conductor}
There is an element $0\neq g\in A$ such that
\[
 gC\subseteq B_1,
 \qquad g^{1/q}C\subseteq B_q\quad(q=p^n).
\]
\end{lemma}

\begin{proof}
The field $E=LK'$ is a finite separable extension of the perfect field
$K'$, and is therefore perfect.  Since it contains $L$ and is contained
in $L_{\mathrm{perf}}$, it equals $L_{\mathrm{perf}}$.
The ring $A'$ is normal: any integral equation and the finitely many
fractions involved in an element of $\operatorname{Frac}(A')$ descend
to some normal ring $A^{1/q}$.

Choose a $K$-basis $b_1,\ldots,b_r$ of $L$ with $b_j\in D$.
Its trace-dual basis $\beta_1,\ldots,\beta_r$ exists by separability.
Since $(A\setminus\{0\})^{-1}D=L$, there is $0\neq g\in A$ with
$g\beta_j\in D$ for every $j$.  The same two bases are trace-dual
after extending scalars to $K'$.

Every $c\in C$ is integral over $A'$, as is $g\beta_j$.
The trace $\operatorname{Tr}_{E/K'}(g\beta_jc)$ is therefore integral
over $A'$ and lies in $K'$, so normality puts it in $A'$.
The trace-dual expansion now reads
\[
 gc=\sum_{j=1}^r
       \operatorname{Tr}_{E/K'}(g\beta_jc)b_j\in A'D=B_1.
\]
Taking $q$-th roots inside the perfect ring $C$ gives the second
inclusion, because $(A'D)^{1/q}=A'D^{1/q}=B_q$.
\end{proof}

We next calculate quotient lengths. Let $H$ be $\mathfrak d$-primary and put
\[
 M_q=B_q/HB_q,\qquad
 N_q=\operatorname{im}(M_q\longrightarrow C/HC).
\]
Finite-stage normalization gives
\begin{equation}\label{eq:perfection-stage-quotient}
 \lambda_\infty(M_q)=q^{-h}\length_D(D/H^{[q]}).
\end{equation}
To see this directly, $M_q$ is the flat base extension of
$D^{1/q}/HD^{1/q}$ from $A^{1/q}$ to $A'$.  The $q$-th power
isomorphism $D^{1/q}\longrightarrow D$ identifies that quotient with
$D/H^{[q]}$, preserving module length.  Length over $A^{1/q}$ equals
length over $D^{1/q}$ because their residue fields agree.

The kernel of $M_q\longrightarrow N_q$ is annihilated by $g^{1/q}$.
In fact, an element in $B_q\cap HC$ is a finite sum of elements of
$H$ times elements of $C$, and Lemma~\ref{lem:perfection-conductor}
puts its product with $g^{1/q}$ in $HB_q$.
For an endomorphism of a module of finite normalized length,
additivity shows that kernel and cokernel have the same normalized
length.  Apply this to multiplication by $g^{1/q}$ on $M_q$ to obtain
\begin{align}
 \lambda_\infty\ker(M_q\longrightarrow N_q)
 &\leq\lambda_\infty(0:_{M_q}g^{1/q})\notag\\
 &=\lambda_\infty(M_q/g^{1/q}M_q)\notag\\
 &=q^{-h}\length_D D/(H^{[q]},g).
 \label{eq:perfection-kernel-bound}
\end{align}
This last quantity is $O(q^{-1})$.  Indeed, choose
$\mathfrak d^a\subseteq H$, and let $t$ be a number of generators of
$H$.  The pigeonhole containment
\[
 H^{t(q-1)+1}\subseteq H^{[q]}
\]
shows that $H^{[q]}$ contains an ordinary power of $\mathfrak d$ whose
exponent is $O(q)$.  If $g$ is a nonunit, the domain hypothesis gives
$\dim D/(g)\leq h-1$, so Hilbert--Samuel growth on $D/(g)$ gives
$\length_D D/(H^{[q]},g)=O(q^{h-1})$.  If $g$ is a unit, this length
is zero.

Combining \eqref{eq:perfection-stage-quotient} and
\eqref{eq:perfection-kernel-bound}, we obtain
\[
 q^{-h}\length_D D/H^{[q]}-O(q^{-1})
 \leq\lambda_\infty(N_q)
 \leq q^{-h}\length_D D/H^{[q]}.
\]
The submodules $N_q$ increase and exhaust $C/HC$.  They are all killed
by $H$ and hence by one fixed power of $\mathfrak m_A$.  The
Hilbert--Kunz limit and Lemma~\ref{lem:normalized-limit} therefore give
\begin{equation}\label{eq:perfection-hk}
 \lambda_\infty(C/HC)=e_{HK}(H,D).
\end{equation}
The quotient lengths are now identified.  To obtain the remaining
parameter assertions of Theorem~\ref{thm:length-inputs}, we first show
that positive parameter Koszul homology has length zero; the Koszul
Euler characteristic will then identify the parameter quotient length
with Hilbert--Samuel multiplicity.  Fix a system of
parameters $z_1,\ldots,z_h$ of $D$.  Flatness at each stage and the
$q$-th power isomorphism give
\begin{equation}\label{eq:perfection-koszul-stage}
 \lambda_\infty H_i(z_1,\ldots,z_h;B_q)
   =q^{-h}\length_D H_i(z_1^q,\ldots,z_h^q;D).
\end{equation}
These are finite lengths: the finite-stage Koszul homology is a
finitely generated module annihilated by a primary parameter ideal.
The complex $K(z_1,\ldots,z_h;D)$ consists of finite free modules in
degrees $0,\ldots,h$, and its homology has finite length.  The
Frobenius homology estimate of Roberts~\cite{Roberts}, in the form
\cite[Theorem~2.4]{Ma}, thus gives
\[
 \lim_{q=p^n\to\infty}
 q^{-h}\length_D H_i(z_1^q,\ldots,z_h^q;D)=0
 \qquad(i>0).
\]
This theorem requires a complete characteristic-$p$ local ring and a
finite-free complex of length equal to its dimension with finite-length
homology, all of which have just been checked.

Filtered colimits are exact for modules, so
\[
 H_i(z_1,\ldots,z_h;C)
   =\varinjlim_q H_i(z_1,\ldots,z_h;B_q).
\]
Each module here is annihilated by $(z_1,\ldots,z_h)$ and therefore
by a fixed power of $\mathfrak m_A$.  The third assertion of
Lemma~\ref{lem:normalized-limit}, applied to
\eqref{eq:perfection-koszul-stage}, proves
\begin{equation}\label{eq:perfection-koszul-zero}
 \lambda_\infty H_i(z_1,\ldots,z_h;C)=0\qquad(i>0).
\end{equation}
The same argument applies to every fixed parameter list
$z_1^a,\ldots,z_h^a$.

Finally, the Koszul Euler characteristic formula for parameter
ideals~\cite[Tag~0AZZ]{Stacks} gives
\[
 \sum_{i=0}^h(-1)^i\length_D
             H_i(z_1^q,\ldots,z_h^q;D)
   =e((z_1^q,\ldots,z_h^q),D)
   =q^h e((z_1,\ldots,z_h),D).
\]
Divide by $q^h$ and use the positive-homology limits.  The degree-zero
term tends to the Hilbert--Kunz multiplicity of the parameter ideal.
Consequently
\[
 e_{HK}((z_1,\ldots,z_h),D)=e((z_1,\ldots,z_h),D).
\]
Together with \eqref{eq:perfection-hk} and the parameter multiplicity
power rule, it gives the quotient formulas in
Theorem~\ref{thm:length-inputs} in characteristic $p$.

\paragraph{Finite-length modules.}
For the Dutta-multiplicity consequence in
Section~\ref{sec:charp}, we also need the same comparison for a
finite-length module, since zeroth homology need not be cyclic.
For a $D$-module $M$, write
\[
 \Phi_D^n(M)=D\otimes_{D,\varphi_D^n}M
\]
for scalar extension along the $n$th Frobenius endomorphism
$\varphi_D^n:D\to D$.

\begin{lemma}\label{lem:perfection-finite-length}
In the characteristic-$p$ setup above, let $M$ be a finite-length
$D$-module.  If the limit on the right exists, then
\[
 \lambda_\infty(C\otimes_D M)
   =\lim_{q=p^n\longrightarrow\infty}
           q^{-h}\length_D\bigl(\Phi_D^n(M)\bigr).
\]
\end{lemma}

\begin{proof}
The assertion is immediate when $M=0$.  Otherwise put
$J=\Ann_D M$, a $\mathfrak d$-primary ideal, and let $b$ be a number
of generators of $M$.  Define
\[
 M_q=B_q\otimes_D M,\qquad
 N_q=\im(M_q\longrightarrow C\otimes_D M).
\]
The module $M_q$ is obtained by flat scalar extension from
$A^{1/q}$ to $A'$ of
$D^{1/q}\otimes_D M$.  Under the $q$th power isomorphism
$D^{1/q}\cong D$, the latter module corresponds to $\Phi_D^n(M)$.
It is a finite-length module over the Noetherian ring $A^{1/q}$,
so finite-stage normalization gives
\begin{equation}\label{eq:perfection-module-stage}
 \lambda_\infty(M_q)=q^{-h}\length_D\bigl(\Phi_D^n(M)\bigr).
\end{equation}
Here the lengths over $A^{1/q}$ and $D^{1/q}$ agree because their
residue fields agree.

By Lemma~\ref{lem:perfection-conductor}, multiplication by $g^{1/q}$
defines a $D$-linear map $C\to B_q$ whose composite
$B_q\hookrightarrow C\to B_q$ is multiplication by $g^{1/q}$.
After tensoring with $M$, this shows that
$K_q=\ker(M_q\to N_q)$ is killed by $g^{1/q}$.  The finite normalized
length in \eqref{eq:perfection-module-stage} permits the same kernel--cokernel
calculation as in \eqref{eq:perfection-kernel-bound}.  Since $M$ is a
quotient of $(D/J)^b$, it gives
\begin{align*}
 \lambda_\infty(K_q)
 &\leq \lambda_\infty(M_q/g^{1/q}M_q)\\
 &\leq b\,\lambda_\infty\bigl(B_q/(J,g^{1/q})B_q\bigr)\\
 &=bq^{-h}\length_D D/(J^{[q]},g)=O(q^{-1}).
\end{align*}
The last estimate is the Hilbert--Samuel growth argument from
\eqref{eq:perfection-kernel-bound}, with $J$ in place of $H$:
$J^{[q]}$ contains a power of
$\mathfrak d$ with exponent $O(q)$, and $\dim D/(g)\leq h-1$ if
$g$ is a nonunit.  If $g$ is a unit, the error is zero.

The submodules $N_q$ increase to $C\otimes_D M$ and are all killed
by $J$, hence by one fixed power of $\mathfrak m_A$.  Additivity,
\eqref{eq:perfection-module-stage}, and
Lemma~\ref{lem:normalized-limit} now give the assertion by taking
the limit.
\end{proof}

\subsection{Mixed characteristic}

Assume now that $D$ has characteristic zero.  Complete $p$ to a system
of parameters $p,t_2,\ldots,t_h$ of $D$.  Cohen structure gives a
module-finite inclusion, with the same residue field,
\[
 A=W(k)[[t_2,\ldots,t_h]]\subseteq D.
\]
It is generically separable because the fraction fields have
characteristic zero.  Let $A_\infty$ be the $p$-adic completion of the
tower above, and set
\[
 C=D^{A_\infty}_{\mathrm{perfd}}
   :=(D\otimes_A A_\infty)_{\mathrm{perfd}}.
\]
The construction uses the perfectoidization of Bhatt--Scholze
\cite[Section~8 and Theorem~10.11]{BhattScholze}; $C$ is the
perfectoid algebra in \cite[Notation~4.0.1]{CLMST}.
Use its normalized length
$\lambda_\infty$, or equivalently restrict scalars to the uncompleted
tower $\Lambda$ when evaluating bounded
$\mathfrak m_A$-power-torsion modules.  Indeed, $\Lambda$ is
$p$-torsion-free, so $\Lambda/p^N\cong A_\infty/p^N$ for every $N$.
A module annihilated by $\mathfrak m_A^N$ is annihilated by $p^N$;
its submodules and finite generating sets are therefore the same
over the two length bases.

For clarity, the exact perfectoid consequences used here are as
follows.  Choose $0\neq g\in A$ such that $A[1/g]\to D[1/g]$ is
finite etale, and write $(g)_{\mathrm{perfd}}$ for the almost ideal in
\cite[Notation~4.0.1]{CLMST}.  Corollary~4.0.4 of that paper says that
the positive Koszul homology on $C$ of every system of parameters of
$A_n\otimes_A D$ is annihilated by $(g)_{\mathrm{perfd}}$.
Taking $n=0$ covers every system of parameters of $D$, not just the
chosen parameters of $A$.
Proposition~4.0.10 identifies the normalized lengths of
$(g)_{\mathrm{perfd}}$-almost-isomorphic
$\mathfrak m_A$-power-torsion modules.  In particular the Koszul
homology just described has length zero.

Definition~4.0.8 and Corollary~4.0.18 of \cite{CLMST} give, for every
$\mathfrak d$-primary ideal $H$,
\[
 \lambda_\infty(C/HC)=e_{\mathrm{perfd}}(H,D)\leq e(H,D),
\]
with equality when $H$ is a parameter ideal.  The equality for
arbitrary parameter ideals is also the $n=0$ case of
Proposition~4.0.14 there.  Applying it to
$(z_1^a,\ldots,z_h^a)$ and using the ordinary parameter multiplicity
power rule yields
\[
 \lambda_\infty\bigl(C/(z_1^a,\ldots,z_h^a)C\bigr)
       =a^h e((z_1,\ldots,z_h),D).
\]
All these statements apply to the same algebra $C$ and the same
normalized length, while $z$ ranges over the systems of parameters
of $D$.

We finish by checking the category and stability assertions in both
characteristics.  A module killed by a power of $\mathfrak d$ is killed
by a power of $\mathfrak m_A$, so the length above is defined on it and
is additive there.  If $H$ is $\mathfrak d$-primary and $H^N M=0$,
choose $r$ with $\mathfrak d^r\subseteq H$; then
$\mathfrak m_A^{rN}M=0$.  Thus a common power of $H$ supplies the common
annihilating power required by Lemma~\ref{lem:normalized-limit}.
The value $e((z),D)$ is positive and finite.  These observations prove
all the remaining assertions of Theorem~\ref{thm:length-inputs}.

\section{Reductions and a finite free complex}\label{sec:reductions}

For maps of positive residue characteristic, our aim is to obtain
a factorization
$R\to T\twoheadrightarrow S$ with $e(T)=e(R)$ and a short free
complex over $T$ supported at its closed point. The first reduction
applies in every characteristic and handles arbitrary dimensions
and residue fields. The construction
that follows also controls the generic local rings of $T$; these
will determine the coefficients in the final multiplicity sums.

\subsection{Completion, components, and residue fields}

We first make reductions which do not impose any restriction on the
characteristic or on the original residue fields.  We use the following
localization theorem for Hilbert--Samuel multiplicity: if $(B,\mathfrak b)$
is an excellent local ring and $\mathfrak q\in\Spec B$ satisfies
\[
  \dim(B/\mathfrak q)+\dim B_{\mathfrak q}=\dim B,
\]
then $e(B_{\mathfrak q})\leq e(B)$.  This is Nagata's localization theorem;
the stated form is recalled in \cite[Theorem~2.1]{Ma}.  It does not require
$B$ to be equidimensional.

\begin{lemma}\label{lem:residue-extension}
For every complete Noetherian local ring $(B,\mathfrak b,k)$ and every
algebraic closure $\overline{k}$ of $k$, there is a complete Noetherian
local ring $B'$ and a flat local map $B\to B'$ such that
\[
  \mathfrak bB'=\mathfrak b_{B'},\qquad
  B'/\mathfrak b_{B'}\cong\overline{k},\qquad e(B')=e(B).
\]
\end{lemma}

\begin{proof}
In equal characteristic, choose a coefficient field and a presentation
$B=k[[X_1,\ldots,X_t]]/I$, and set
\[
  B'=\overline{k}[[X_1,\ldots,X_t]]/
             I\overline{k}[[X_1,\ldots,X_t]].
\]
The map between the power series rings is flat.  For example, this
follows from the local flatness criterion modulo $(X_1,\ldots,X_t)$:
the variables form regular sequences in both rings, and the resulting
map of coefficient fields is flat.  The quotient map is therefore flat
as well and has the asserted maximal ideal and residue field.

In the remaining case the residue characteristic is $p>0$, and the
Cohen Structure Theorem~\cite[Tag~032A]{Stacks} gives a presentation
$B=V[[X_1,\ldots,X_t]]/I$, where $V$ is a complete discrete
valuation ring with uniformizer $p$ and residue field $k$.  We explain
why $V$ admits a complete discrete valuation extension $V'$ with the
same uniformizer and residue field $\overline{k}$, even when $k$ is
imperfect.  Given a simple algebraic residue extension with irreducible
minimal polynomial $\overline f$, lift $\overline f$ to a monic
polynomial $f\in V[X]$.  The algebra $V[X]/(f)$ is finite free over $V$,
and its quotient by $p$ is the desired residue field.  Every maximal
ideal lies over $(p)$, so this algebra is local with maximal ideal
generated by $p$.  The element $p$ is a nonzerodivisor, and the algebra
has dimension one.  It is consequently a discrete valuation ring.
This argument does not require $\overline f$ to be separable.

Repeat this construction along a tower of simple algebraic extensions
whose union is $\overline{k}$.  At limit stages take unions.  Every
nonzero element of such a union is $p^a$ times a unit for an integer
$a\geq0$; hence every nonzero ideal is generated by an element of least
valuation.  The union is again a discrete valuation ring with
uniformizer $p$.  Completing the final union gives the required $V'$.
The extension $V\to V'$ is flat because $V'$ is torsion-free over $V$.
Now use
\[
  B'=V'[[X_1,\ldots,X_t]]/IV'[[X_1,\ldots,X_t]].
\]
Flatness of the power series extension follows from the same local
flatness criterion, this time with coefficient map $V\to V'$.

Finally, for any flat local map with extended maximal ideal, a
composition series of a finite-length module remains a composition
series after tensoring: each residue-field factor becomes the target
residue field.  In particular
\[
  \length_{B'}(B'/\mathfrak b^N B')
       =\length_B(B/\mathfrak b^N)\quad(N\geq1).
\]
The two Hilbert--Samuel functions, and therefore their degrees and
multiplicities, agree.
\end{proof}

\begin{lemma}\label{lem:reduction}
To prove $e(R)\leq e(S)$ for all flat local maps of nonzero Noetherian
local rings, it suffices to prove it when $R$ and $S$ are complete,
$R$ is a domain, $\dim R=\dim S>0$, and the residue field of $S$ is
algebraically closed.  The reduction preserves equal characteristic
$p$ when both original rings have characteristic $p$.
\end{lemma}

\begin{proof}
Completion preserves dimension and Hilbert--Samuel multiplicity, and
a flat local map induces a flat map on completions.  We can thus first
assume that $R$ and $S$ are complete.  Put $d=\dim R$ and
$f=\dim(S/\mathfrak mS)$.  The dimension formula for a flat local map
gives $\dim S=d+f$.  Choose $Q$ minimal over $\mathfrak mS$ with
$\dim(S/Q)=f$.  Going-down gives $\dim S_Q\geq d$, whereas
\[
  \dim S\geq\dim(S/Q)+\dim S_Q\geq f+d=\dim S.
\]
All these inequalities are equalities.  Since $S$ is excellent,
Nagata's theorem gives $e(S_Q)\leq e(S)$.  The map $R\to S_Q$ is flat
local, its closed fiber has dimension zero, and $\dim S_Q=d$.
Replacing $S_Q$ by its completion preserves these properties.  Thus a
proof for this new target proves the original inequality.

Choose a finite prime filtration of the $R$-module $R$, with factors
$R/\mathfrak p_j$, allowing repetitions.  Flatness gives a filtration
of $S$ with factors $S/\mathfrak p_jS$.  The flat local map
\[
  R/\mathfrak p_j\longrightarrow S/\mathfrak p_jS
\]
has the same zero-dimensional closed fiber, so its source and target
have equal dimension.  The additivity formula for multiplicity yields
\[
 e(R)=\sum_{\dim(R/\mathfrak p_j)=d}e(R/\mathfrak p_j),\qquad
 e(S)=\sum_{\dim(R/\mathfrak p_j)=d}e(S/\mathfrak p_jS).
\]
It therefore suffices to prove the inequality for the indicated domain
sources.  Their quotients and the corresponding targets remain
complete.  If $d=0$, each such source is a field and has multiplicity
one, while a nonzero zero-dimensional local target has multiplicity
equal to its positive length.  This case is settled.

For $d>0$, apply Lemma~\ref{lem:residue-extension} to the target.
The resulting composite remains flat local and has the same target
dimension and multiplicity.  This proves the reduction.  The argument
also appears in \cite[Lemma~2.2]{Ma}.
\end{proof}

\subsection{A factorization with a perfect kernel}

The passage from a flat local map to a short free complex follows the
construction used by Iyengar--Ma--Walker
\cite[Proposition~8.3]{IMW2022}. We construct one auxiliary ring and one
complex for both cases of positive residue characteristic, while also
controlling the local rings needed to sum the component estimates. A
\emph{perfect ideal} $J$ of a Noetherian local ring $T$ means an ideal
for which $T/J$ has finite projective dimension and
\[
  \operatorname{pd}_T(T/J)=\operatorname{grade}_T J.
\]

\begin{proposition}\label{prop:setup}
Let $(R,\mathfrak m)\to(S,\mathfrak n)$ be a flat local map of complete
Noetherian local rings.  Suppose $R$ is a domain,
$\dim R=\dim S=d>0$, the residue characteristic is $p>0$, and the
residue field $L$ of $S$ is algebraically closed.  There is a factorization
\[
  R\longrightarrow (T,\mathfrak t)\longrightarrow S=T/J
\]
with the following properties.  Write
$b=\dim(T/\mathfrak mT)$ and $h=d+b$.
\begin{enumerate}
\item[(i)] The ring $T$ is complete, its residue field is $L$, and
      $R\to T$ is flat with regular closed fiber.  Moreover
      $\dim T=h$ and $e(T)=e(R)$.
\item[(ii)] Every irreducible component of $T$ has dimension $h$.
      The ideal $J$ is perfect of grade $b$, and $S$ has a minimal
      finite free $T$-resolution $P_\bullet$ in degrees $0,\ldots,b$
      with $P_0=T$.
\item[(iii)] Choose lifts $x_1,\ldots,x_d\in\mathfrak t$ of a minimal
      reduction of $\mathfrak n$.  Then
      \[
         F_\bullet=P_\bullet\otimes_T K_\bullet(x_1,\ldots,x_d;T)
      \]
      is a finite free complex in degrees $0,\ldots,h$, has $F_0=T$,
      and is contractible locally on
      $\Spec T\setminus\{\mathfrak t\}$.  Its differential entries
      lie in $\mathfrak t$, and its ranks have alternating sum zero.
\item[(iv)] If $Q\in\Spec T$ contracts to zero in $R$, then $T_Q$ is
      a complete intersection.  In mixed characteristic, $T_Q$ is
      regular.  Every prime of $T$ which is minimal over $J$ contracts
      to zero in $R$.  In mixed characteristic $T$ is generically
      reduced; a prime $Q$ minimal over $J$ contains a unique minimal
      prime of $T$.
\end{enumerate}
\end{proposition}

Parts (i)--(iii) preserve the source multiplicity and produce the complex
to which Theorem~\ref{thm:complex} will apply. Part (iv) supplies the
characteristic-dependent comparisons at primes $Q$ minimal over $J$.
In characteristic $p$, its complete-intersection assertion will identify
$e_{\mathrm{HK}}(JT_Q,T_Q)$ with $\length_{T_Q}(T_Q/JT_Q)$ in
Lemma~\ref{lem:ma-hk}. In mixed characteristic, regularity puts each
such prime on exactly one component of $T$, so the later summation
counts it only once.

\begin{proof}
\emph{Constructing $T$.}
Let $k$ be the residue field of $R$.  Choose a complete regular
normalization $A\subset R$ with residue field $k$, such that $R$ is
finite over $A$.  In equal characteristic, take a coefficient field
$k\subset R$ and a system of parameters $a_1,\ldots,a_d$, giving
$A=k[[a_1,\ldots,a_d]]$.  In mixed characteristic, choose a coefficient
Cohen ring $V\subset R$ and a system of parameters
$p,a_2,\ldots,a_d$, giving $A=V[[a_2,\ldots,a_d]]$.  The Cohen
Structure Theorem and completeness give finiteness in both cases;
the maps from these power series rings are injective by dimension.
No generic separability is required here.

Apply the Cohen Factorization Theorem of
Avramov--Foxby--Herzog~\cite{AFH} to $A\to S$:
\[
  A\longrightarrow (U,\mathfrak u)\twoheadrightarrow S,
\]
where $U$ is complete, $A\to U$ is flat local, and its closed fiber is
regular; see \cite[Theorem~3.1]{Ma}.  Regularity ascends along this map
because both $A$ and the closed fiber are regular
\cite[Tag~031E]{Stacks}. Thus $U$ is a
regular local ring, with residue field $L$.

Set $T=R\otimes_A U$.  This is finite over $U$, and multiplication of
the two given maps to $S$ induces a surjection $T\to S$.  Its closed
fiber over $U$ is
\[
  T/\mathfrak uT=(R/\mathfrak m_A R)\otimes_k L.
\]
The ring $R/\mathfrak m_A R$ is local Artinian with residue field $k$.
Its radical is nilpotent, so the displayed fiber is local Artinian
with residue field $L$.  A finite algebra over a complete local ring
is a product of complete local rings; this fiber shows that there is
only one factor.  Hence $T$ is complete local with residue field $L$.
Base change shows that $R\to T$ is flat and
\[
  T/\mathfrak mT\cong U/\mathfrak m_A U.
\]
Its closed fiber is therefore regular of dimension $b$, and the flat
dimension formula gives $\dim T=d+b=h=\dim U$.

\begin{figure}[ht]
\centering
\[
\begin{array}{ccccc}
 A&\longrightarrow&U&&\\
 \big\downarrow&&\big\downarrow&&\\
 R&\longrightarrow&T=R\otimes_A U&\twoheadrightarrow&S .
\end{array}
\]
\caption{The construction of the auxiliary factorization. The square
is a tensor-product square; $R$ is finite over $A$, and $A\to U$
is flat local with regular closed fiber. The surjection to $S$
comes from its given $R$- and $U$-algebra structures.}
\label{fig:factorization}
\end{figure}
Figure~\ref{fig:factorization} records the two operations in the
construction: the finite normalization $A\subset R$ and the Cohen
factorization of $A\to S$.

\smallskip
\noindent\emph{Preserving multiplicity.}
To check the multiplicity assertion, lift a regular system of
parameters of $T/\mathfrak mT$ to $w_1,\ldots,w_b\in\mathfrak t$.
Lifting regular sequences under a flat local map
\cite[Tag~00MG]{Stacks} shows that $w$ is a
$T$-regular sequence and $T/(w)$ is flat over $R$.  Completeness
defines the evaluation map
\[
   R[[W_1,\ldots,W_b]]\longrightarrow T,\qquad W_i\longmapsto w_i.
\]
The local flatness criterion modulo the variables
\cite[Tag~00ML]{Stacks} makes this map flat.
Its extended maximal ideal is $\mathfrak t$, because the images of
$w$ generate the maximal ideal of the closed fiber.  As in the last
paragraph of Lemma~\ref{lem:residue-extension}, the Hilbert--Samuel
functions are unchanged under a flat local map with extended maximal
ideal.  Adjoining formal variables preserves multiplicity, as follows
also from the associated graded polynomial extension.  Therefore
\[
  e(T)=e(R[[W_1,\ldots,W_b]])=e(R).
\]

\smallskip
\noindent\emph{The resolution and the perfect kernel.}
Start with a minimal free $T$-resolution of $S$, possibly infinite.
Each syzygy is flat over $R$:
inductively, the kernel of a surjection between $R$-flat modules with
flat quotient is flat.  Tensoring with $k$ over $R$ consequently
preserves exactness and yields a minimal free resolution of
$S/\mathfrak mS$ over the regular local ring $T/\mathfrak mT$.
That ring has global dimension $b$.  Minimality, or Nakayama's Lemma,
therefore forces the original resolution to vanish in degrees greater
than $b$.  Thus $\operatorname{pd}_T S\leq b$.

The closed fiber $S/\mathfrak mS$ is zero-dimensional, so the image of
$J$ in $T/\mathfrak mT$ is primary for its maximal ideal.  Choose in
this image a system of $b$ parameters and lift them to elements of
$J$.  Since the fiber is regular, this is a regular sequence on the
fiber; flat regular-sequence lifting makes its lifts a $T$-regular
sequence.  Consequently
\[
 b\leq\operatorname{grade}_T J
   \leq\operatorname{pd}_T S\leq b.
\]
This proves perfection and identifies the length of the minimal
resolution.  The assertions include $b=0$, when the regular sequence
is empty and the resolution has only its degree-zero term.

\smallskip
\noindent\emph{Components and the supported complex.}
Because $R$ is a finite torsion-free $A$-module, it embeds in a finite
free $A$-module.  Flatness gives an embedding of $T$ in a finite free
$U$-module.  The contractions to $U$ of the associated primes of $T$
are thus zero.  In particular every minimal prime $P$ of $T$ contracts
to zero, and $T/P$ is finite integral over $U$.  It follows that
$\dim(T/P)=\dim U=h$.

The field $L$ is infinite, so a minimal reduction of $\mathfrak n$
generated by $d$ elements exists.  Lift its generators to $x$ in $T$.
The ideal $J+(x)$ is $\mathfrak t$-primary.  At a prime outside its
vanishing set, either the localized resolution resolves zero and is
split exact, or one of the $x_i$ is a unit and its Koszul complex is
contractible.  Hence $F_\bullet$ is contractible off $\mathfrak t$.
Minimality of the resolution and $x_i\in\mathfrak t$ put all its
differential entries in $\mathfrak t$.  Its degree-zero term is $T$,
and
\[
  \sum_i(-1)^i\rank F_i
    =\left(\sum_i(-1)^i\rank P_i\right)(1-1)^d=0,
\]
where $d>0$ is used.

\smallskip
\noindent\emph{Local rings over the generic point of $R$.}
Finally put $K=\operatorname{Frac}(A)$ and
$K_R=\operatorname{Frac}(R)$.  Localizing at $A\setminus\{0\}$ gives
\[
 T\otimes_A K
   \cong (U\otimes_A K)\otimes_K K_R.
\]
Here $U\otimes_A K$ is regular, and $K_R/K$ is a finite field
extension.  Express this field extension as a finite tower of simple
algebraic extensions.  After base change each step is defined by a
monic polynomial, which is a nonzerodivisor even if the intermediate
ring is not reduced.  The resulting local rings are quotients of
regular local rings by regular sequences, hence complete
intersections.  Every $Q$ contracting to zero in $R$ belongs to this
localization, proving the first assertion of (iv).  In mixed
characteristic the fields have characteristic zero, so $K_R/K$ is
separable.  The displayed extension is then finite etale, and its
local rings are regular.

Every minimal prime of $S$ contracts to a minimal prime of $R$ by
flat going-down; that prime is zero.  This applies to the primes of
$T$ minimal over $J$.  Moreover every minimal prime of $T$ avoids
$A\setminus\{0\}$ by the preceding torsion-freeness argument, so in
mixed characteristic its localization is a zero-dimensional regular
local ring, namely a field.  Thus $T$ is generically reduced.  At a
prime $Q$ minimal over $J$, the regular local ring $T_Q$ is a domain,
which also shows that $Q$ contains exactly one minimal prime of $T$.
\end{proof}

\begin{remark}\label{rem:setup-dimensions}
Since $T$ is complete and equidimensional, it is catenary and
\[
  \dim T_Q+\dim(T/Q)=h\qquad(Q\in\Spec T).
\]
In particular a prime $Q$ containing $J$ with $\dim(T/Q)=d$ is minimal
over $J$ and has $\dim T_Q=b$.  In mixed characteristic its
localization is regular by Proposition~\ref{prop:setup}.
\end{remark}

\section{The characteristic-\texorpdfstring{$p$}{p} argument}\label{sec:charp}

In characteristic $p$, Frobenius makes the differential entries of
the complex from Proposition~\ref{prop:setup} arbitrarily deep.
We apply Theorem~\ref{thm:complex} separately to each component of
$T$ and add the resulting inequalities. The following
Hilbert--Kunz identity identifies that sum with $e(S)$.

\begin{lemma}[Ma's Hilbert--Kunz identity]\label{lem:ma-hk}
Let $(T,\mathfrak t)$ be a complete local ring of characteristic
$p>0$ with perfect residue field, let $J$ be a perfect ideal, and put
$S=T/J$.  For any system of parameters $x_1,\ldots,x_d$ of $S$,
\[
 e_{\mathrm{HK}}(J+(x),T)
   =\sum_{\substack{Q\text{ minimal over }J\\\dim(T/Q)=d}}
        e((x),T/Q)\,e_{\mathrm{HK}}(JT_Q,T_Q).
\]
If each $T_Q$ occurring in this sum is a complete intersection, then
\[
       e_{\mathrm{HK}}(J+(x),T)=e((x),S).
\]
\end{lemma}

\begin{proof}
The first assertion is \cite[Lemma~4.2]{Ma}. For the second, we use
the following complete-intersection consequence from
\cite[Section~2.3 and Lemma~4.2]{Ma}: a primary ideal of finite
projective dimension in a local complete intersection of
characteristic $p$ has Hilbert--Kunz multiplicity equal to its
colength. Since $J$ is perfect, $JT_Q$ is such an ideal in $T_Q$.
Consequently
\[
 e_{\mathrm{HK}}(JT_Q,T_Q)=\length_{T_Q}(T_Q/JT_Q).
\]
The Hilbert--Samuel associativity formula gives the conclusion.  No
assumption that $x$ extend to a parameter system of $T$ is involved.
\end{proof}

\begin{theorem}\label{thm:positive-characteristic}
For every flat local map $(R,\mathfrak m)\to(S,\mathfrak n)$ of
nonzero Noetherian local rings of characteristic $p>0$, one has
$e(R)\leq e(S)$.
\end{theorem}

\begin{proof}
Lemma~\ref{lem:reduction} reduces the problem to
complete rings of the same positive dimension, with $R$ a domain and
the residue field of $S$ algebraically closed.  Use
Proposition~\ref{prop:setup} to obtain $T,J,P_\bullet,F_\bullet$ and
the integers $d,b,h=d+b$.  Set
\[
       H=J+(x_1,\ldots,x_d)T.
\]
This is a $\mathfrak t$-primary ideal.

Fix a minimal prime $P$ of $T$ and put $D=T/P$, with maximal ideal
$\mathfrak d$.  This is a complete local domain of dimension $h$ and
has the same algebraically closed residue field.  Choose a minimal
reduction $(z_1,\ldots,z_h)$ of $\mathfrak d$.  There is an integer
$c\geq0$ such that
\[
  \mathfrak d^n=(z)^{n-c}\mathfrak d^c
      \subseteq(z)^{n-c}\qquad(n\geq c).
\]
Take the $D$-algebra $C$ and its normalized length $\lambda$ from
Theorem~\ref{thm:length-inputs}.  For $I=(z)C$ the length hypotheses
of Theorem~\ref{thm:complex} hold with
\[
       \mu=e((z),D)=e(D).
\]

For a power $q=p^n$, let $F_\bullet(q)$ be the complex obtained by
raising every matrix entry of the differentials of $F_\bullet$ to
its $q$th power.  This is scalar extension by the $n$th Frobenius
endomorphism of $T$ and has the same ranks as $F_\bullet$.  Its
differential entries lie in $\mathfrak t^{[q]}$.  Their images in
$C$ therefore lie in
\[
     \mathfrak d^{[q]}C\subseteq\mathfrak d^qC
                          \subseteq I^{q-c}.
\]
Frobenius preserves the identities defining a contracting homotopy
on each localization where $F_\bullet$ is contractible.  Subsequent
scalar extension preserves those identities as well.  Since $(z)$
is $\mathfrak d$-primary, the ideals $I$ and $\mathfrak tC$ have
the same radical.  Consequently
$F_\bullet(q)\otimes_T C$ is contractible locally off $V(I)$.

The first differential of the resolution $P_\bullet$ has entries
generating $J$, and $P_0=T$.  Thus
\[
 H_0(F_\bullet(q)\otimes_T C)=C/H^{[q]}C.
\]
Apply Theorem~\ref{thm:complex} with $s=q-c$, for $q>c$.  Its
degree-zero rank is one, and the ranks of all its terms are fixed
as $q$ varies.  The Hilbert--Kunz quotient formula from
Theorem~\ref{thm:length-inputs} and Frobenius scaling give
\begin{align*}
 q^h e_{\mathrm{HK}}(HD,D)
   &=e_{\mathrm{HK}}((HD)^{[q]},D)\\
   &=\lambda(C/H^{[q]}C)\\
   &\geq e(D)(q-c)^h(1-o(1)).
\end{align*}
Dividing by $q^h$ and letting $q$ tend to infinity proves
\[
             e_{\mathrm{HK}}(HD,D)\geq e(D).
\]

Every component of $T$ has dimension $h$ by
Proposition~\ref{prop:setup}.  The associativity formulas for
Hilbert--Kunz~\cite[Theorem~3.14]{HunekeHK} and Hilbert--Samuel
multiplicities have the same
generic length coefficients, so
\begin{align*}
 e_{\mathrm{HK}}(H,T)
   &=\sum_{P\in\operatorname{Min}(T)}
       \length_{T_P}(T_P)\,
       e_{\mathrm{HK}}(H(T/P),T/P)\\
   &\geq\sum_{P\in\operatorname{Min}(T)}
       \length_{T_P}(T_P)\,e(T/P)
     =e(T).
\end{align*}
In particular this summation does not assume that $T$ is reduced.

We now verify the hypotheses of Lemma~\ref{lem:ma-hk}.  The ring $T$
is complete of characteristic $p$, its residue field is algebraically
closed and therefore perfect, and $J$ is perfect.  For every minimal
prime $Q$ of $J$ with $\dim(T/Q)=d$, Proposition~\ref{prop:setup}
shows that $T_Q$ is a complete intersection.  The chosen $x$ is a
system of parameters of $S=T/J$, and its ideal is a reduction of
$\mathfrak n$.  Hence
\[
  e(S)=e((x),S)=e_{\mathrm{HK}}(H,T)
          \geq e(T)=e(R),
\]
as required.
\end{proof}

\subsection{Dutta multiplicity of short complexes}

The zeroth homology in the following consequence need not be cyclic.
The finite-length-module comparison in
Lemma~\ref{lem:perfection-finite-length} therefore takes the place of
the Hilbert--Kunz quotient formula used above.
We use the following terminology of Iyengar--Ma--Walker
\cite[Sections~2.5--2.6]{IMW2022}.

\begin{definition}\label{def:short-complex}
Let $(R,\mathfrak m)$ be a Noetherian local ring of characteristic
$p>0$ and dimension $d$.  A \emph{short complex} over $R$ is a
complex of finite free modules
\[
 F=(0\longrightarrow F_d\longrightarrow\cdots
                 \longrightarrow F_0\longrightarrow0)
\]
whose homology modules have finite length and for which $H_0(F)\ne0$.
For any bounded complex $G$ with finite-length homology, write
\[
 \chi_R(G)=\sum_i(-1)^i\length_R H_i(G).
\]
When $G$ is finite free, let $\Phi_R^nG$ denote scalar extension
along the $n$th Frobenius endomorphism.  Its \emph{Dutta multiplicity}
is
\[
 \chi_\infty^R(G)=
 \lim_{n\longrightarrow\infty}
       p^{-nd}\chi_R(\Phi_R^nG).
\]
\end{definition}

The limit exists; see \cite[Section~4.17]{IMW2022}.  For a short
complex over a complete characteristic-$p$ local ring, the Frobenius
homology theorem of Roberts, as recalled in
\cite[Theorem~2.4]{Ma}, states that
\[
 \lim_{n\longrightarrow\infty}
 p^{-nd}\length_R H_i(\Phi_R^nF)=0\qquad(i>0).
\]
Consequently
\begin{equation}\label{eq:dutta-h0}
 \chi_\infty^R(F)=
 \lim_{n\longrightarrow\infty}
 p^{-nd}\length_R H_0(\Phi_R^nF).
\end{equation}
This holds for arbitrary short complexes; it does not require their
zeroth homology to be cyclic.

Iyengar, Ma, and Walker conjectured the lower bound in the next
corollary for short complexes over all complete local rings
\cite[Conjecture~8.1]{IMW2022}.  Here we prove the bound for complete
domains of positive characteristic.

\begin{corollary}[Dutta multiplicity over positive-characteristic domains]
\label{cor:dutta-domain}
Let $(R,\mathfrak m)$ be a complete Noetherian local domain of
characteristic $p>0$, and let $F$ be a short complex over $R$.
Then
\[
                    \chi_\infty^R(F)\geq e(R).
\]
\end{corollary}

\begin{proof}
Put $d=\dim R$.  If $d=0$, then $R$ is a field and $F_0$ is a
nonzero free module, so $\chi_\infty^R(F)=\rank_R F_0\geq1=e(R)$.
Assume $d>0$.  Splitting off contractible free summands makes $F$
minimal without changing its homology or Dutta multiplicity.
Write $b_i=\rank_R F_i$; in particular $b_0\geq1$.

First suppose that the residue field is algebraically closed.  Choose
a minimal reduction $I=(z_1,\ldots,z_d)$ of $\mathfrak m$.  As in the
proof of Theorem~\ref{thm:positive-characteristic}, there is a fixed
integer $c\geq0$ with
\[
                 \mathfrak m^q\subseteq I^{q-c}\qquad(q\geq c).
\]
Use the characteristic-$p$ algebra $C=R_{\mathrm{perf}}$ and normalized
length $\lambda$ from Theorem~\ref{thm:length-inputs}.  For any short
complex $G$ over $R$, right exactness gives
\[
 C\otimes_R H_0(G)=H_0(C\otimes_R G),\qquad
 \Phi_R^n(H_0(G))=H_0(\Phi_R^nG).
\]
Lemma~\ref{lem:perfection-finite-length} and
\eqref{eq:dutta-h0} therefore give
\[
                  \lambda H_0(C\otimes_R G)=\chi_\infty^R(G).
\]

Let $q=p^n$.  The entries of the differentials of $\Phi_R^nF$ lie
in $\mathfrak m^{[q]}\subseteq\mathfrak m^q\subseteq I^{q-c}$.
Frobenius preserves a contracting homotopy at every localization
where $F$ is contractible.  Thus $\Phi_R^nF$ again has finite-length
homology; it is minimal and has nonzero zeroth homology, so it is
short.  Scalar extension to $C$ preserves the same contracting
homotopies; since $I$ is $\mathfrak m$-primary, the extended complex
is contractible locally off $V(IC)$.  The defining Dutta limit, with
its index shifted by $n$, now gives
\[
 \lambda H_0(C\otimes_R\Phi_R^nF)
   =\chi_\infty^R(\Phi_R^nF)=q^d\chi_\infty^R(F).
\]
Moreover $\sum_i(-1)^ib_i=0$, since $F$ becomes exact over the
fraction field of $R$.  Apply Theorem~\ref{thm:complex} to
$C\otimes_R\Phi_R^nF$ with parameter ideal $IC$, depth $s=q-c$,
and $\mu=e(I,R)=e(R)$.  For $q>c$ it gives
\[
 q^d\chi_\infty^R(F)
 \geq e(R)(q-c)^d
       \left(b_0-\frac{A_d}{q-c}\sum_{i=0}^d b_i\right).
\]
After division by $q^d$, the right side tends to $b_0e(R)$.
Since $b_0\geq1$, this proves the corollary in this case.

For an arbitrary residue field, Lemma~\ref{lem:residue-extension}
gives a complete flat local extension
$(R,\mathfrak m)\to(S,\mathfrak n)$ with
$\mathfrak n=\mathfrak mS$, algebraically closed residue field, and
$\dim S=d$, $e(S)=e(R)$.  Put $G=S\otimes_R F$.  This is again a
minimal short complex.  Frobenius scalar extension commutes with
this base change, and finite-length module lengths are preserved,
so
\[
             \chi_\infty^S(G)=\chi_\infty^R(F);
\]
this invariance is also recorded in \cite[Section~4.17]{IMW2022}.
The ring $S$ need not be a domain, so we pass to its components by
a prime filtration rather than apply the preceding case directly.

Choose a finite prime filtration of the $S$-module $S$ with factors
$D_j=S/P_j$, and put $h_j=\dim D_j$.  Additivity of Euler
characteristic for the bounded free complex $\Phi_S^nG$ gives
\[
 \chi_S(\Phi_S^nG)
       =\sum_j\chi_S(D_j\otimes_S\Phi_S^nG).
\]
The complexes in the sum have finite-length homology because the
contracting homotopies off $\mathfrak n$ persist after tensoring.
The Hochster--Huneke Frobenius homology estimate, in the form
\cite[Theorem~5.2]{BhattHochsterMa}, gives constants $C_j$ such that
\[
 \length_S H_i(D_j\otimes_S\Phi_S^nG)
       \leq C_jp^{n\min\{h_j,d-i\}}
       \qquad(0\leq i\leq d).
\]
Indeed, its complex has length $d$ and finite-length homology, and
its finite coefficient module $D_j$ has dimension $h_j$.
When $h_j<d$, every homology length in this summand is
$O(p^{nh_j})$, so its contribution divided by $p^{nd}$ tends to
zero.

When $h_j=d$, the ring $D_j$ is a complete local domain with the
same algebraically closed residue field.  The complex
$D_j\otimes_SG$ is short over $D_j$: its homology is supported at
the maximal ideal because the localized contracting homotopies
persist, and its zeroth homology is nonzero because its matrices
remain minimal and $b_0\geq1$.  Inspecting differential matrices
gives
\[
 D_j\otimes_S\Phi_S^nG
       \cong\Phi_{D_j}^n(D_j\otimes_SG).
\]
Lengths of its homology over $S$ and $D_j$ agree.  Taking limits,
applying the already proved domain case to these factors, and using
additivity of top-dimensional Hilbert--Samuel multiplicity along the
prime filtration, we obtain
\[
 \chi_\infty^R(F)=\chi_\infty^S(G)
   =\sum_{h_j=d}\chi_\infty^{D_j}(D_j\otimes_SG)
   \geq\sum_{h_j=d}e(D_j)=e(S)=e(R).
\]
This completes the proof.
\end{proof}

\section{Mixed characteristic}\label{sec:mixed}

We work one irreducible component at a time. Finite integral extensions
make the differential entries deep enough to apply the uniform complex
estimate. An upper bound for the resulting normalized length is obtained
from a prime filtration, and is independent of those extensions after a
separate parameter-power limit. The following two multiplicity estimates
supply that upper bound. Whenever an
ideal is used on a quotient ring, its image in that quotient is understood.

\begin{lemma}\label{lem:mixed-quotient}
Let $(A,\mathfrak a)$ be a Noetherian local ring of positive dimension
$t$, let $H$ be $\mathfrak a$-primary, and let $f\in H$.  If
$\dim A/(f)=t-1$, then
\[
  e(H,A)\leq e(H,A/(f)).
\]
\end{lemma}

\begin{proof}
Multiplication by $f$ induces a map $A/H^{j-1}\longrightarrow A/H^j$
whose cokernel is $A/(H^j,f)$.  Consequently
\[
 \length_A(A/(H^j,f))
 \geq \length_A(A/H^j)-\length_A(A/H^{j-1}).
\]
For large $j$ these are Hilbert--Samuel polynomials.  The polynomial
on the right has leading coefficient $e(H,A)/(t-1)!$, and the one
on the left has leading coefficient $e(H,A/(f))/(t-1)!$.
Comparing them proves the assertion, including the case $t=1$.
\end{proof}

\begin{lemma}\label{lem:mixed-asymptotic}
Let $(B,\mathfrak b)$ be a complete Noetherian local domain of
dimension $h$, and let $1\leq d\leq h$. Suppose that an ideal
$L\subseteq\mathfrak b$ and elements $x_1,\ldots,x_d\in\mathfrak b$ satisfy
$\sqrt{L+(x_1,\ldots,x_d)}=\mathfrak b$.  Put
\[
 H_N=L+(x_1^N,\ldots,x_d^N),\qquad N\geq1.
\]
The following assertions hold.
\begin{enumerate}
\item If $\dim B/L<d$, then $e(H_N,B)=O(N^{d-1})$.
\item If $L=\mathfrak q$ is prime, $\dim B/\mathfrak q=d$, and
$B_{\mathfrak q}$ is regular, then
\[
 e(H_N,B)\leq N^d e((x_1,\ldots,x_d),B/\mathfrak q)
                 +O(N^{d-1}).
\]
\end{enumerate}
All constants concern the fixed data $B,L,x$ and are independent of $N$.
\end{lemma}

\begin{proof}
Choose $a\geq1$ with $\mathfrak b^a\subseteq L+(x_1,\ldots,x_d)$,
and set $r_N=d(N-1)+1$.  The pigeonhole principle for monomials gives
\begin{equation}\label{eq:mixed-power-containment}
 \mathfrak b^{a r_N}
 \subseteq \bigl(L+(x_1,\ldots,x_d)\bigr)^{r_N}
 \subseteq L+(x_1^N,\ldots,x_d^N)=H_N.
\end{equation}
In particular, on any fixed nonzero quotient $A$ of $B$ of dimension
$t\leq d-1$, monotonicity of multiplicity under inclusion of primary
ideals gives
\begin{equation}\label{eq:mixed-small-quotient}
 e(H_N,A)\leq e(\mathfrak b^{a r_N},A)
       =(a r_N)^t e(\mathfrak b,A)=O(N^{d-1}).
\end{equation}
This also applies when $t=0$.

To use this estimate, we will pass from $B$ to fixed quotients of
smaller dimension.  Quotienting by an element of $L$ that lowers the
dimension by one can only increase the multiplicity of $H_N$, by
Lemma~\ref{lem:mixed-quotient}.  In the first case we will reach
dimension $d-1$.  In the second case we will first reach dimension
$d$, choosing the quotient so that its localization at $\mathfrak q$
is the residue field of $B_{\mathfrak q}$.  The component
$B/\mathfrak q$ will therefore have coefficient one in the
associativity formula, while every other component of that quotient
will admit one more dimension-lowering quotient.

The ring $B$ is catenary and equidimensional.  Thus, for every prime
$\mathfrak p\subset B$, one has
$\dim B/\mathfrak p=h-\operatorname{ht}\mathfrak p$.
We will also use the following consequence.  A list of $j$ partial
parameters in $B$ has all its minimal primes of height $j$: its
quotient has dimension $h-j$, which gives the lower bound on these
heights, and the generalized principal ideal theorem gives the
upper bound.  In particular its quotient is equidimensional.

For the first assertion, select $h-d+1$ partial parameters
$v_1,\ldots,v_{h-d+1}$ in $L$.  Here is the required prime avoidance
argument.  After $j<h-d+1$ choices, a minimal prime of the preceding
ideal has height $j$ and quotient dimension $h-j\geq d$.  It cannot
contain $L$, because $\dim B/L<d$.  Hence the next element can be
chosen in $L$ outside all these finitely many primes.  Each choice
lowers the dimension by one.  Applying
Lemma~\ref{lem:mixed-quotient} successively, with $H_N$ fixed, gives
\[
 e(H_N,B)\leq e(H_N,B/(v_1,\ldots,v_{h-d+1})).
\]
The quotient on the right is fixed and has dimension $d-1$, so
\eqref{eq:mixed-small-quotient} proves the first assertion.

For the second assertion put $c=h-d=\operatorname{ht}\mathfrak q$.
We choose $v_1,\ldots,v_c\in\mathfrak q$ that are partial parameters
in $B$ and generate the maximal ideal of the regular local ring
$B_{\mathfrak q}$.  To justify both conditions simultaneously,
suppose that $v_1,\ldots,v_j$ have been chosen, with $j<c$, and that
their images in
\[
 \mathfrak qB_{\mathfrak q}/(\mathfrak qB_{\mathfrak q})^2
\]
are linearly independent over $\kappa(\mathfrak q)$.  Let $E_j$ be
the contraction to $B$ of
\[
 (v_1,\ldots,v_j)B_{\mathfrak q}
                 +(\mathfrak qB_{\mathfrak q})^2.
\]
Since this cotangent space has dimension $c$, the ideal
$\mathfrak q$ is not contained in $E_j$.  Nor is it contained in
any minimal prime of $(v_1,\ldots,v_j)$: those primes have height
$j<c$.  Prime avoidance, with the single possibly nonprime ideal
$E_j$, therefore supplies an element $v_{j+1}$ in $\mathfrak q$
outside their union.  This preserves both conditions.  At the end,
Nakayama's Lemma gives
\[
 (v_1,\ldots,v_c)B_{\mathfrak q}=\mathfrak qB_{\mathfrak q}.
\]
For $c=0$ this is the empty list.

Set $A=B/(v_1,\ldots,v_c)$, which is equidimensional of dimension
$d$.  Repeated application of Lemma~\ref{lem:mixed-quotient}, followed
by associativity of multiplicity, gives
\begin{equation}\label{eq:mixed-cut-associativity}
 e(H_N,B)\leq e(H_N,A)
   =\sum_{\mathfrak p\in\operatorname{Min}(A)}
      \length_{B_{\mathfrak p}}(A_{\mathfrak p})\,
      e(H_N,B/\mathfrak p),
\end{equation}
where primes of $A$ are identified with primes of $B$.
The prime $\mathfrak q$ is one of these minimal primes, and its
coefficient is
\[
 \length_{B_{\mathfrak q}}(A_{\mathfrak q})
 =\length_{B_{\mathfrak q}}
      (B_{\mathfrak q}/\mathfrak qB_{\mathfrak q})=1.
\]
On this component $H_N$ becomes $(x_1^N,\ldots,x_d^N)$.
The parameter-ideal power formula therefore makes its contribution
exactly $N^d e((x_1,\ldots,x_d),B/\mathfrak q)$.

For any other prime $\mathfrak p$ in
\eqref{eq:mixed-cut-associativity}, the two distinct primes
$\mathfrak p$ and $\mathfrak q$ have the same height $c$ and are
incomparable.  Choose $w_{\mathfrak p}\in\mathfrak q\setminus
\mathfrak p$.  Its image is nonzero in the complete local domain
$B/\mathfrak p$ of dimension $d$, and its quotient has dimension
$d-1$.  It belongs to the image of every $H_N$.  Thus
Lemma~\ref{lem:mixed-quotient} and
\eqref{eq:mixed-small-quotient} give
\[
 e(H_N,B/\mathfrak p)
 \leq e(H_N,B/(\mathfrak p,w_{\mathfrak p}))=O(N^{d-1}).
\]
The number of these components and their coefficients in
\eqref{eq:mixed-cut-associativity} are fixed and finite.  Summing
proves the second assertion.
\end{proof}

\begin{theorem}\label{thm:mixed-characteristic}
Assume the notation and conclusions of Proposition~\ref{prop:setup},
and suppose that the source domain $R$ has mixed characteristic.
Then
\[
 e(S)\geq e(T)=e(R).
\]
\end{theorem}

\begin{proof}
Write $\mathfrak t$ for the maximal ideal of $T$, and recall that
$\dim T=h$, $\dim S=d>0$, $S=T/J$, and
\[
 F=P\otimes_T K(x_1,\ldots,x_d;T),
\]
where $P$ is a finite free resolution of $S$, of length at most
$h-d$.  The complex $F$ has length at most $h$, has $F_0=T$, and is
contractible away from the closed point.  Its differentials have
entries in $\mathfrak t$, and its ranks have alternating sum zero.
The images of $x_1,\ldots,x_d$ form a minimal reduction of the maximal
ideal of $S$; in particular $\sqrt{J+(x_1,\ldots,x_d)}=\mathfrak t$.
The residue field, denoted by $k$, is algebraically closed.

Fix a minimal prime $P_0$ of $T$ and put
$B=T/P_0$, with maximal ideal $\mathfrak b$.  By
Proposition~\ref{prop:setup}, $B$ is a complete local domain of mixed
characteristic, has dimension $h$, and has residue field $k$.
Choose a parameter ideal $(u_1,\ldots,u_h)$ that reduces
$\mathfrak b$.

\medskip
\noindent\emph{Finite extensions that make the differentials deep.}
Fix an integer $q>h$ and choose roots $z_i=u_i^{1/q}$ in an algebraic
closure of $\operatorname{Frac}(B)$.  The ring
$B'=B[z_1,\ldots,z_h]$ is a finite domain extension of $B$.
It is complete local, has residue field $k$, and has dimension $h$.
The ideal $(z_1,\ldots,z_h)$ is primary for its maximal ideal.
Since $\mathfrak b$ is integral over $(u_1,\ldots,u_h)$, every
element of $\mathfrak b$ is integral over $(z_1,\ldots,z_h)^qB'$.

Choose $a\geq1$ with $p^a\in(z_1,\ldots,z_h)B'$, where $p$ is
the residue characteristic.  Apply Heitmann's integral-extension
form of the Brian\c{c}on--Skoda Theorem~\cite{Heitmann}, as stated in
\cite[Proposition~2.8]{HeitmannMa} to the generating list
\[
 p^a,z_1,\ldots,z_h
\]
of the ideal $I=(z_1,\ldots,z_h)B'$.  This list has $n=h+1$
members and satisfies the theorem's hypothesis
$p\in\sqrt{(p^a,z_1)B'}$.  With $k'=q-h-1\geq0$, its conclusion
sends each element integral over $I^{n+k'}=I^q$ into
$I^{k'+1}=I^{q-h}$ after an integral extension.
Here redundant generators are allowed by the theorem's statement.
Apply it to a finite generating set of $\mathfrak b$.  Taking a
prime above zero in each extension if necessary gives integral
domain extensions, which can be embedded in a common algebraic
closure of $\operatorname{Frac}(B)$.  Adjoin the finitely many
integral coefficients witnessing the resulting memberships.
We obtain a domain $D$ finite over $B$ such that
\begin{equation}\label{eq:mixed-deep-entries}
 \mathfrak bD\subseteq(z_1,\ldots,z_h)^{q-h}D.
\end{equation}
We may replace $D$ by its normalization: it remains finite by
excellence, and the inclusion persists.  A finite algebra over a
complete local ring is a product of complete local rings; since
$D$ is a domain, there is just one factor.  Thus $D$ is complete
local of dimension $h$, and its residue field is $k$ because it
is finite over the algebraically closed field $k$.

Put $r=\rank_B D$.  For every $\mathfrak b$-primary ideal $H$,
the finite-module multiplicity formula and equality of residue
fields give
\begin{equation}\label{eq:mixed-rank-formula}
 e(HD,D)=r e(H,B).
\end{equation}
Indeed, a $\operatorname{Frac}(B)$-basis of
$D\otimes_B\operatorname{Frac}(B)$ can be
cleared of denominators to give an injection $B^r\longrightarrow D$
with torsion cokernel.  This cokernel has dimension less than $h$,
so it does not contribute to multiplicity.  Moreover, lengths of
finite-length $D$-modules are the same over $B$ and over $D$ because
the residue fields agree.  Applying
\eqref{eq:mixed-rank-formula} to $(u_1,\ldots,u_h)$ and using the
parameter-ideal power formula gives
\begin{equation}\label{eq:mixed-mu-normalization}
 q^h e((z_1,\ldots,z_h),D)
 =e((u_1,\ldots,u_h),D)=r e(B).
\end{equation}

\medskip
\noindent\emph{A lower bound for the normalized length.}
Use the algebra $C$ and additive length $\lambda$ supplied for $D$
by Theorem~\ref{thm:length-inputs}, and put
$\mu=e((z_1,\ldots,z_h),D)$.  Extend $F$ along
$T\longrightarrow B\longrightarrow D\longrightarrow C$.
Its entries belong to $(z_1,\ldots,z_h)^{q-h}C$ by
\eqref{eq:mixed-deep-entries}.  Local contractions survive this
base change.  The ideals $\mathfrak tD$ and $(z_1,\ldots,z_h)D$
have the same radical, so the required support is also preserved.
Theorem~\ref{thm:complex}, with $s=q-h$ and zeroth rank one,
therefore gives a constant $K$ depending only on $h$ and the fixed
ranks of $F$ such that
\[
 \lambda(C/(J,x_1,\ldots,x_d)C)
 \geq \mu(q-h)^h\left(1-\frac{K}{q-h}\right)
\]
for all sufficiently large $q$.  Together with
\eqref{eq:mixed-mu-normalization}, this reads
\begin{equation}\label{eq:mixed-lower-bound}
 \frac{\lambda(C/(J,x_1,\ldots,x_d)C)}{r}
 \geq e(B)\left(1-\frac hq\right)^h
                  \left(1-\frac{K}{q-h}\right).
\end{equation}
In particular, $K$ is independent of the choices of $D$, $C$, and
$\lambda$ as $q$ varies.

\medskip
\noindent\emph{The power rule for the quotient lengths.}
For now keep $q$, $D$, $C$, and $\lambda$ fixed.  For every list of
positive integers $a_1,\ldots,a_d$, consider
\[
 F(a)=P_C\otimes_C K(x_1^{a_1},\ldots,x_d^{a_d};C).
\]
Before base change this complex is contractible away from the
closed point: outside $V(J)$ the resolution $P$ is split exact,
and outside $V(x_1,\ldots,x_d)$ the Koszul complex is contractible.
Its length is at most $h$.  Lemma~\ref{lem:acyclicity} therefore
shows that its positive homology has $\lambda$-length zero and
that its zeroth homology has finite $\lambda$-length.  Thus
\[
 \chi_\lambda(F(a))
 :=\sum_i(-1)^i\lambda(H_i(F(a)))
 =\lambda(C/(J,x_1^{a_1},\ldots,x_d^{a_d})C).
\]

To check the power rule without any finiteness assumption on
partial Koszul complexes, fix all exponents but $a_i$, and let
$E$ denote $P_C$ tensored with the remaining Koszul factors.
For positive integers $v,w$, the composable scalar maps
$x_i^v,x_i^w$ on $E$ give the mapping-cone triangle
\[
 \operatorname{Cone}(x_i^v:E\to E)
 \longrightarrow\operatorname{Cone}(x_i^{v+w}:E\to E)
 \longrightarrow\operatorname{Cone}(x_i^w:E\to E)
 \longrightarrow\operatorname{Cone}(x_i^v:E\to E)[1].
\]
Every cone includes a positive power of each $x_j$, so all three
have finite homology lengths by the preceding argument.  Additivity
of $\lambda$ in their long exact homology sequence gives additivity
of Euler characteristic.  Hence that Euler characteristic is
linear in each positive exponent.  Applying this to all $d$
exponents proves
\begin{equation}\label{eq:mixed-length-power-rule}
 \lambda(C/(J,x_1^N,\ldots,x_d^N)C)
   =N^d\lambda(C/(J,x_1,\ldots,x_d)C)
 \qquad(N\geq1).
\end{equation}
No assertion that the homology of $E$ itself has finite length was
needed.

\medskip
\noindent\emph{An upper bound from a prime filtration.}
Fix a finite prime filtration
\[
 0=S_0\subset S_1\subset\cdots\subset S_m=S,
 \qquad S_j/S_{j-1}\simeq T/Q_j.
\]
All $Q_j$ contain $J$; the list records repetitions.  Set
\[
 M_N=C/(x_1^N,\ldots,x_d^N)C,
 \qquad H_{j,N}=Q_jB+(x_1^N,\ldots,x_d^N)B.
\]
Tensoring each filtration sequence with $M_N$ gives a right exact
sequence
\[
 S_{j-1}\otimes_T M_N\longrightarrow S_j\otimes_T M_N
                   \longrightarrow C/H_{j,N}C\longrightarrow0.
\]
Every module here is killed by $J+(x_1^N,\ldots,x_d^N)$, whose
extension to $D$ is primary for its maximal ideal.  Their lengths
are finite: for example, each $S_j\otimes_T M_N$ is a quotient of
a finite direct sum of $C/(J,x_1^N,\ldots,x_d^N)C$.
The image of the first arrow is a quotient of its source, so
additivity and nonnegativity of $\lambda$ yield
\[
 \lambda(S_j\otimes_T M_N)
 \leq\lambda(S_{j-1}\otimes_T M_N)+\lambda(C/H_{j,N}C).
\]
Summing these inequalities, applying the primary-ideal bound from
Theorem~\ref{thm:length-inputs}, and then applying
\eqref{eq:mixed-rank-formula}, we obtain
\begin{equation}\label{eq:mixed-filtration-bound}
 \lambda(C/(J,x_1^N,\ldots,x_d^N)C)
 \leq\sum_{j=1}^m\lambda(C/H_{j,N}C)
 \leq r\sum_{j=1}^m e(H_{j,N},B).
\end{equation}
Thus this use of the prime filtration requires only right
exactness of tensor product, not flatness of $C$ over $T$.

We apply Lemma~\ref{lem:mixed-asymptotic} to each term of the last
sum.  Since $Q_j\supseteq J$, one has $\dim T/Q_j\leq d$.
If $\dim T/Q_j<d$, then also $\dim B/Q_jB<d$.  If $P_0$ is not
contained in $Q_j$, its image in the domain $T/Q_j$ is nonzero,
and hence
\[
 \dim B/Q_jB=\dim T/(P_0+Q_j)<\dim T/Q_j\leq d.
\]
In either case the first assertion of the Lemma gives
$e(H_{j,N},B)=O(N^{d-1})$.

The remaining case is $\dim T/Q_j=d$ and $P_0\subseteq Q_j$.
Such a $Q_j$ is minimal over $J$: a strictly smaller prime over
$J$ would give a quotient of dimension greater than $d$.
The generic regularity in Proposition~\ref{prop:setup} gives that
$T_{Q_j}$ is regular.  It follows that $Q_j$ contains exactly one
minimal prime of $T$, namely $P_0$, and
\[
 B_{Q_jB}=T_{Q_j}.
\]
Indeed the localization of $P_0$ is the unique minimal prime of
the regular local domain $T_{Q_j}$ and is therefore zero.
The second assertion of Lemma~\ref{lem:mixed-asymptotic} now gives
\begin{equation}\label{eq:mixed-component-asymptotic}
 e(H_{j,N},B)
 \leq N^d e((x_1,\ldots,x_d),T/Q_j)+O(N^{d-1}).
\end{equation}

\medskip
\noindent\emph{The two limits and the component sum.}
Divide \eqref{eq:mixed-filtration-bound} by $rN^d$ and use
\eqref{eq:mixed-length-power-rule}.  Let $N$ tend to infinity
while keeping the chosen $q,D,C,\lambda$ fixed.  The preceding
asymptotic estimates imply
\begin{equation}\label{eq:mixed-uniform-upper-bound}
 \frac{\lambda(C/(J,x_1,\ldots,x_d)C)}{r}
 \leq
 \sum_{\substack{1\leq j\leq m\,:\ \dim T/Q_j=d\\P_0\subseteq Q_j}}
       e((x_1,\ldots,x_d),T/Q_j).
\end{equation}
The right side does not involve any auxiliary extension or length
algebra.  Combine this bound with
\eqref{eq:mixed-lower-bound} and then let $q$ tend to infinity.
Uniformity of $K$ gives
\begin{equation}\label{eq:mixed-one-component}
 e(T/P_0)\leq
 \sum_{\substack{1\leq j\leq m\,:\ \dim T/Q_j=d\\P_0\subseteq Q_j}}
       e((x_1,\ldots,x_d),T/Q_j).
\end{equation}
There is no exchange of limits: the $N$-limit was taken separately
for each fixed choice of the data belonging to $q$.

Finally sum \eqref{eq:mixed-one-component} over the minimal primes
$P_0$ of $T$.  The ring $T$ is equidimensional and generically
reduced by Proposition~\ref{prop:setup}; its generic local lengths
are therefore all one, and associativity gives
\[
 e(T)=\sum_{P_0\in\operatorname{Min}(T)}e(T/P_0).
\]
Every $Q_j$ of dimension $d$ belongs to exactly one of these
components, because $T_{Q_j}$ is regular.  Thus summation counts
each occurrence of $Q_j$ in the filtration exactly once.  The
number of occurrences of a fixed such prime $Q$ is
$\length_{T_Q}(S_Q)$, as seen by localizing the filtration at $Q$.
The multiplicity formula for $S$ consequently yields
\[
 e(T)
 \leq\sum_{\substack{1\leq j\leq m\,:\ \dim T/Q_j=d}}
           e((x_1,\ldots,x_d),T/Q_j)
 =e((x_1,\ldots,x_d),S)=e(S).
\]
The equality $e(T)=e(R)$ from Proposition~\ref{prop:setup}
finishes the proof.
\end{proof}

\section{Residue characteristic zero}\label{sec:charzero}

We deduce the residue-characteristic-zero case from
Theorem~\ref{thm:positive-characteristic} by encoding a hypothetical
strict counterexample in finitely many polynomial equations. We follow
the reduction-and-approximation strategy of \cite[Section~5]{Ma},
beginning with the finite free reduction of \cite[Lemma~5.1]{Ma}.
For the subsequent transfer, it is enough to preserve a lower bound for
the source multiplicity and an upper bound for the target multiplicity. This
allows us to use a presentation matrix and finitely many initial
relations, without keeping track of the homology of a resolution.

\subsection{A finite free map with a common coefficient field}

\begin{lemma}\label{lem:charzero-finite-free}
Let $(R,\mathfrak m)\to(S,\mathfrak n)$ be a flat local map between
complete Noetherian local rings of equal dimension and residue
characteristic zero.  There is a factorization
\[
 R\longrightarrow R'\longrightarrow S
\]
such that $R'$ is complete local, $e(R')=e(R)$, and $R'\to S$ is finite
free.  Moreover, $R'$ and $S$ have the same coefficient field and the
same residue field.
\end{lemma}

\begin{proof}
Choose a coefficient field $k\subset R$, and identify it with its image
in $S$.  Put $L=S/\mathfrak n$.  A coefficient field of $S$ containing
this copy of $k$ can be chosen as follows.  Lift a transcendence basis
of $L/k$ to $S$.  These lifts generate a copy of the corresponding
rational function field, because a polynomial with nonzero residue is
a unit whenever it occurs as a denominator.  The remaining algebraic
extension of residue fields is separable.  Hensel's lemma lifts each
simple algebraic root, and a maximal compatible choice, obtained by
Zorn's lemma, lifts all of $L$.  We henceforth identify this coefficient
field with $L$.

Choose a presentation $R=k[[X_1,\ldots,X_a]]/I$ and set
\[
 R'=L[[X_1,\ldots,X_a]]/I L[[X_1,\ldots,X_a]].
\]
The compatible maps from $R$ and $L$ to $S$ give a local map $R'\to S$.
Extension of coefficient fields on power series rings is flat; for
example, the local criterion of flatness applies to the regular
sequence of variables and the field extension on their quotients.
Thus $R\to R'$ is flat, its maximal ideal is $\mathfrak m'=\mathfrak mR'$,
and its residue field is $L$.  Filtering by powers of the maximal
ideal gives
\[
 \gr_{\mathfrak m'}R'
   \cong (\gr_{\mathfrak m}R)\otimes_k L.
\]
In particular, $\dim R'=\dim R$ and $e(R')=e(R)$.

The flat dimension formula for the original map shows that
$S/\mathfrak mS$ is Artinian.  Since
$\mathfrak m'S=\mathfrak mS$, this quotient is a finite-dimensional
$L$-vector space.  Lift an $L$-basis to $s_1,\ldots,s_q\in S$.
Successively expressing a remainder modulo $\mathfrak m'S$, then
modulo $(\mathfrak m'S)^2$, and so on, expresses every element of $S$
as $\sum_i a_i s_i$ with $a_i\in R'$.  Indeed, the coefficient
corrections at step $j$ lie in $(\mathfrak m')^j$, so they converge in
$R'$, and the remainders tend to zero in $S$: the
$\mathfrak m'S$-adic and $\mathfrak n$-adic topologies agree.  Hence
$S$ is finite over $R'$.

Let $P_\bullet$ be a free resolution of $k$ over $R$.  Flatness of
$R'$ over $R$ makes $P_\bullet\otimes_R R'$ a free resolution of $L$
over $R'$.  Therefore
\[
 \Tor^{R'}_1(L,S)
   = H_1(P_\bullet\otimes_R S)
   = \Tor^R_1(k,S)=0.
\]
A finite module over a Noetherian local ring with this vanishing is
free.  Thus $R'\to S$ is finite free, necessarily of positive rank.
\end{proof}

\subsection{Encoding the multiplicity bounds}

\begin{lemma}\label{lem:charzero-certificate}
Let $(R_0,\mathfrak m_0)\to(S_0,\mathfrak n_0)$ be a finite free
local map of complete local rings of dimension $d>0$, with a common
infinite coefficient field $L$ and residue field $L$.  Put
$r=e(R_0)$.  There is a finite system
\[
 F(t_1,\ldots,t_d,Y_1,\ldots,Y_N)=0
\]
of polynomial equations over $\ZZ$ with a solution in
$L[[t_1,\ldots,t_d]]$, having the following property.  For every field
$K$ and every solution in $A=K[[t_1,\ldots,t_d]]$, the solution
determines a finite free local map
\[
 (R,\mathfrak a)\longrightarrow(S,\mathfrak b)
\]
of nonzero complete local rings of dimension $d$ and residue field
$K$, for which
\[
 e(R)\ge r,\qquad e(S)\le e(S_0).
\]
\end{lemma}

\begin{proof}
Choose a minimal reduction $t_1,\ldots,t_d$ of $\mathfrak m_0$.
The induced map
\[
 A_0=L[[t_1,\ldots,t_d]]\longrightarrow R_0
\]
is finite, by completeness and the parameter property.  It is
injective: a nonzero kernel in the regular local domain $A_0$ would
force the dimension of its finite algebra $R_0$ to be less than $d$.
Since the residue fields agree, the finite-module multiplicity formula
and the reduction property give
\begin{equation}\label{eq:charzero-original-rank}
 r=e(R_0)=e((t),R_0)=\rank_{A_0}R_0.
\end{equation}
Here no domain or Cohen--Macaulay hypothesis on $R_0$ is needed.
For completeness, a finite module of rank $r$ over $A_0$ contains a
copy of $A_0^r$ with torsion cokernel; the cokernel has dimension less
than $d$, so additivity of multiplicity gives the last equality.

We describe the finite system.  Throughout the construction, an
equality in a presented module means a coordinate equality with an
additional vector witnessing membership in its relation submodule.
Thus, for a matrix $U$, the assertion $w\in\operatorname{im}U$ is
encoded by $w=Uc$ with a new vector of unknowns $c$.  The number of
coordinates and assertions below is always fixed and finite.

The system records the algebra structures, freeness, and locality.
The numerical conditions use the source's generic rank for
a lower multiplicity bound and the dimensions of the target's associated
graded pieces for an upper bound.

\smallskip
\noindent\emph{The source algebra and its generic rank.}
Fix a presentation
\[
 A_0^v\xrightarrow{U_0}A_0^u\longrightarrow R_0\longrightarrow0.
\]
The entries of a matrix $U$ of this same size are unknowns.  On
$A^u$, specify a bilinear multiplication by vectors $c_{ij}\in A^u$
giving the products of standard basis vectors, and specify a unit
vector $\epsilon\in A^u$.  Impose the following identities modulo
$\operatorname{im}U$: commutativity and associativity on basis
vectors, the two unit identities on basis vectors, and the assertion
that a relation column multiplied by any basis vector is again in
$\operatorname{im}U$.  These identities give a unital commutative
$A$-algebra structure on
\[
 R=\operatorname{coker}U.
\]
They are polynomial equations in the entries of $U$, the
multiplication table, $\epsilon$, and the membership witnesses.

Impose also the vanishing of all $(u-r+1)$-minors of $U$.  If there
are no minors of that size, this condition is empty.  Over the
fraction field of $A$, it gives
\begin{equation}\label{eq:charzero-rank-bound}
 \rank_A R=u-\rank U\ge r>0.
\end{equation}
These equations are satisfied by the original presentation in view of
\eqref{eq:charzero-original-rank}. A solution may give $U$ smaller rank,
but this only increases the rank of its cokernel. To turn that rank
bound into a multiplicity bound, we will retain the source reduction
$(t)$ as well.

\smallskip
\noindent\emph{The target as a free source-module.}
Let $n>0$ be the original free rank of $S_0$ over $R_0$.
Define the underlying module of the new target to be $S=R^n$.
Give it an $R$-bilinear multiplication by specifying
\[
 \eta_i\eta_j=\sum_{h=1}^n d_{ijh}\eta_h,
 \qquad d_{ijh}\in R,
\]
and a unit $\delta\in R^n$.  Represent these coefficients and this
unit by vectors over $A$.  Impose commutativity, associativity, and
the unit identities on the $R$-basis $\eta_1,\ldots,\eta_n$.
The underlying $A$-module is presented by the block diagonal matrix
with $n$ copies of $U$, so these are again finitely many polynomial
equations, using the source multiplication table and relation
witnesses.  The structural map is $a\mapsto a\delta$.
Consequently, in every solution, $S$ is a unital commutative
$R$-algebra which is free of rank $n$ as an $R$-module.

\smallskip
\noindent\emph{The local ideals and the source reduction.}
Choose finite generating lists for $\mathfrak m_0$ and
$\mathfrak n_0$ that include the respective images of $t_1,\ldots,t_d$.
Represent corresponding lists in $R$ and $S$ by vectors of
unknowns, with their first $d$ entries defined to be
$t_i\epsilon$ and $t_i\delta$.  Denote the generated ideals by
$\mathfrak a$ and $\mathfrak b$.  Fix exponents $N_R,N_S$ and $j$
for which the original rings satisfy the following conditions, and
impose them on the new data:
\begin{equation}\label{eq:charzero-ideal-conditions}
 \mathfrak a^{N_R}\subseteq(t)R,\qquad
 \mathfrak b^{N_S}\subseteq(t)S,\qquad
 \mathfrak a^{j+1}\subseteq(t)\mathfrak a^j.
\end{equation}
Such exponents exist because $(t)$ is primary for both original
maximal ideals and reduces $\mathfrak m_0$.
For each standard $A$-module generator $e_i$ of $R$, impose an
expression
\[
 e_i=c_i\epsilon+\sum_\nu a_\nu w_{i\nu},
 \qquad c_i\in A,\quad w_{i\nu}\in R,
\]
where $a_\nu$ is the source ideal list.  Impose the analogous
expressions for the $un$ standard $A$-module generators of $S$,
using its unit and its ideal list.  Thus $R/\mathfrak a$ and
$S/\mathfrak b$ are spanned by their units as $A$-modules.
All these conditions reduce to identities on finite lists of
generators and monomials.  Expanding products with the multiplication
tables, and representing each coefficient by a vector, encodes them
as polynomial equations.

We first check their consequences, before imposing the target
Hilbert-function conditions.  By \eqref{eq:charzero-rank-bound}, $R$
is a nonzero finite $A$-algebra, and $S=R^n$ is also nonzero and
finite over $A$.  In a finite algebra over the local ring $A$, every
maximal ideal contains the extended ideal $(t)$.  The first two
conditions in \eqref{eq:charzero-ideal-conditions} therefore show
that every maximal ideal of $R$ contains $\mathfrak a$, and every
maximal ideal of $S$ contains $\mathfrak b$.  In particular, both
ideals are proper.  Since $R/\mathfrak a$ is spanned by the unit and
annihilated by $(t)$, it is a nonzero quotient of $K$, hence is $K$.
Thus $\mathfrak a$ is the unique maximal ideal of $R$.  The same
argument gives $S/\mathfrak b=K$ and uniqueness of $\mathfrak b$.

The map $R\to S$ is local.  Indeed, an element of $\mathfrak a$
is nilpotent modulo $(t)R$, so its image belongs to
$\sqrt{(t)S}=\mathfrak b$; contraction of $\mathfrak b$ is then
$\mathfrak a$.  Positive generic rank implies that $R$ is faithful
as an $A$-module, so $A\to R$ is injective.  The map $R\to S$ is
also injective: if $a\delta=0$, multiplication by $a$ on $S$ is
zero, whereas $S=R^n$ is faithful.  Finiteness now gives
$\dim R=\dim S=d$.  Both rings are complete, since they are finite
over complete $A$ and their maximal-ideal topologies agree with their
$(t)$-adic topologies by \eqref{eq:charzero-ideal-conditions}.

Because $(t)\subseteq\mathfrak a$, the last inclusion in
\eqref{eq:charzero-ideal-conditions} is equality.  It makes $(t)$
a reduction of $\mathfrak a$, whence
\begin{equation}\label{eq:charzero-source-bound}
 e(R)=e((t),R)=\rank_A R\ge r.
\end{equation}
The rank equality has the same proof as
\eqref{eq:charzero-original-rank}. Thus the reduction relation converts
the lower rank bound into the required source multiplicity bound.
Neither a nonvanishing minor nor an additional condition on the length
of the residue quotient is required.

\smallskip
\noindent\emph{An upper bound for the target Hilbert function.}
Write $y_1,\ldots,y_q$ for the chosen original target ideal list.
There is a graded surjection
\[
 L[Z_1,\ldots,Z_q]\longrightarrow\gr_{\mathfrak n_0}S_0,
 \qquad Z_i\longmapsto y_i+\mathfrak n_0^2.
\]
Let $I$ be its homogeneous kernel.  Fix a degree-compatible monomial
order.  Choose finitely many monic homogeneous polynomials
$f_1,\ldots,f_b\in I$ whose leading monomials generate the initial
ideal $J=\operatorname{in}(I)$.  Write
\[
 f_\ell=Z^{\alpha_\ell}
       +\sum_{\beta\in B_\ell}c_{\ell\beta}Z^\beta,
 \qquad |\beta|=|\alpha_\ell|=a_\ell,
 \quad Z^\beta<Z^{\alpha_\ell}.
\]
Here $B_\ell$ is a fixed finite set.  Replace the coefficients
$c_{\ell\beta}$ by unknowns $C_{\ell\beta}\in A$ and require, on
the new target ideal list $b_1,\ldots,b_q$, that
\begin{equation}\label{eq:charzero-leading-relations}
 b^{\alpha_\ell}
  +\sum_{\beta\in B_\ell}C_{\ell\beta}b^\beta
       \in\mathfrak b^{a_\ell+1}.
\end{equation}
Membership is expressed using the finitely many degree
$a_\ell+1$ monomials in the ideal list and coefficient vectors in
$S$.  These are again polynomial equations, and the original data
provide a solution.

For any transferred solution, reduce the coefficients
$C_{\ell\beta}$ modulo $(t)$.  Equation
\eqref{eq:charzero-leading-relations} gives a homogeneous relation in
the kernel of $K[Z]\to\gr_{\mathfrak b}S$ with the same leading
monomial $Z^{\alpha_\ell}$.  The leading coefficient is still $1$.
The higher $(t)$-adic terms of the other coefficients do not contribute:
their products with a degree-$a_\ell$ monomial lie in
$\mathfrak b^{a_\ell+1}$.  It follows that the new initial ideal
contains the monomial ideal $J$, interpreted over $K$. The transferred
target may have additional initial relations; this containment is the
direction needed for an upper bound. Standard monomials outside $J$
thus span each new graded piece, whereas they
form a basis of the corresponding original graded piece.  Consequently
\[
 \dim_K\mathfrak b^a/\mathfrak b^{a+1}
 \le \dim_L\mathfrak n_0^a/\mathfrak n_0^{a+1}
 \quad(a\ge0).
\]
Summation and the equality of dimensions give
\begin{equation}\label{eq:charzero-target-bound}
 e(S)\le e(S_0).
\end{equation}

Finally, every coefficient in the foregoing presentations, tables,
lists, and membership witnesses is an unknown; only the sizes,
exponents, monomials, and the distinguished coefficients $1$ have
been fixed.  Accordingly the entire system has coefficients in
$\ZZ[t_1,\ldots,t_d]$.  It is finite, the original data solve it in
$A_0$, and \eqref{eq:charzero-source-bound} and
\eqref{eq:charzero-target-bound} prove the assertion.
\end{proof}

\subsection{Specialization to positive characteristic}

The preceding encoding reduces the problem to transferring existence
of a solution of a finite polynomial system. Artin approximation supplies
an algebraic solution whose coefficients can be specialized.

\begin{lemma}\label{lem:charzero-specialization}
Let $L$ be a field of characteristic zero.  A finite polynomial system
over $\ZZ[t_1,\ldots,t_d]$ which has a solution in $L[[t_1,\ldots,t_d]]$
has a solution in $K[[t_1,\ldots,t_d]]$ for some algebraically closed
field $K$ of positive characteristic.
\end{lemma}

\begin{proof}
Put $B=(L[t_1,\ldots,t_d]_{(t)})^h$, whose completion is $L[[t]]$.
Artin approximation \cite[Theorem~1.10]{Artin} gives a solution in
$B$.  The henselization is the filtered direct limit of pointed
etale neighborhoods of the origin.  Thus the finitely many solution
entries and identities descend to one such neighborhood with an
$L$-rational point over $t=0$; this is also the explicit conclusion
of \cite[Corollary~2.1]{Artin}.

Shrink about that point.  An etale presentation then has the form
\[
 E=\bigl(L[t,z_1,\ldots,z_m]/(G_1,\ldots,G_m)\bigr)_{g\Delta},
 \qquad
 \Delta=\det\left(\frac{\partial G_i}{\partial z_j}\right),
\]
where the distinguished point is $(t,z)=(0,c)$, with
$c\in L^m$, $G_i(0,c)=0$, and $g(0,c)\Delta(0,c)\ne0$.
The solution entries are elements of $E$.  Choose polynomial
representatives with denominators powers of $g\Delta$.  For each
of the finitely many solution identities, clear denominators and
choose a polynomial expression in the ideal $(G_1,\ldots,G_m)$
which witnesses that identity.  One may multiply by a further power
of $g\Delta$ before doing so.

Let $W\subset L$ be the $\ZZ$-algebra generated by the following
finite set: the coefficients of $G_i$ and $g$, the coordinates of
$c$, the coefficients of the representatives and ideal-membership
witnesses just chosen, and the inverse of $g(0,c)\Delta(0,c)$.
This is a nonzero finitely generated $\ZZ$-algebra of characteristic
zero.  The same presentation defines an etale $W[t]$-algebra
\[
 E_W=\bigl(W[t,z]/(G_1,\ldots,G_m)\bigr)_{g\Delta}.
\]
Evaluation at $(0,c)$ gives a $W$-valued point.  All solution
identities hold in $E_W$, by the chosen witnesses.

Choose a maximal ideal of $W$ and let $k_0$ be its residue field.
This field is finite.  Indeed, it is a field finitely generated as a
$\ZZ$-algebra.  If it had characteristic zero, Zariski's lemma over
$\mathbb Q$ would make it a number field.  Finitely many ring
generators of a number field are contained in its ring of integers
with one nonzero integer inverted, which cannot contain the inverse
of every rational prime.  This is impossible for a field of
characteristic zero.  In positive characteristic, Zariski's lemma
makes $k_0$ a finite extension of its finite prime field.

Let $K$ be an algebraic closure of $k_0$.  The distinguished point
specializes to a $K$-rational point of $E_W\otimes_W K$ over $t=0$.
The inverted Jacobian and denominator remain nonzero there.
The etale lifting property, or multivariable Hensel's lemma, lifts
this point to a $K[t]$-algebra map
\[
 E_W\otimes_W K\longrightarrow K[[t]].
\]
Explicitly, the point lifts uniquely over each successive quotient
$K[[t]]/(t)^a$, since the Jacobian is invertible, and the compatible
lifts give the displayed map by completeness.  Applying it to the
solution entries preserves all the polynomial identities and proves
the assertion.
\end{proof}

\begin{theorem}\label{thm:charzero}
For every flat local homomorphism of nonzero Noetherian local rings
$(R,\mathfrak m)\to(S,\mathfrak n)$ whose residue fields have
characteristic zero, one has $e(R)\le e(S)$.
\end{theorem}

\begin{proof}
By Lemma~\ref{lem:reduction}, it suffices to consider complete rings
of the same dimension with algebraically closed target residue field.
Lemma~\ref{lem:charzero-finite-free} then permits us to assume that
the map is finite free and that both rings have the same
algebraically closed coefficient field $L$.

If the common dimension is zero, equality of residue fields and
freeness of rank $n\ge1$ give
$e(S)=\length_S S=n\length_R R\ge e(R)$.
Suppose instead that $d>0$ and that $e(R)>e(S)$.
Lemma~\ref{lem:charzero-certificate} produces a finite polynomial
system with a solution over $L[[t_1,\ldots,t_d]]$.
Lemma~\ref{lem:charzero-specialization} gives a solution over
$K[[t_1,\ldots,t_d]]$ for an algebraically closed field $K$ of
positive characteristic.  The resulting finite free local map
$\widetilde R\to\widetilde S$ satisfies
\[
 e(\widetilde R)\ge e(R)>e(S)\ge e(\widetilde S),
\]
contradicting Theorem~\ref{thm:positive-characteristic}.
\end{proof}


\section{Proof of the main theorem}\label{sec:assembly}
\begin{proof}[Proof of Theorem~\ref{thm:main}]
The reductions of Section~\ref{sec:reductions} reduce the assertion to
a complete, equal-dimensional flat local map with domain source and
algebraically closed target residue field; the dimension-zero case is
immediate. A source of positive characteristic is covered by
Theorem~\ref{thm:positive-characteristic}, a source of mixed
characteristic by Theorem~\ref{thm:mixed-characteristic}, and a source
of residue characteristic zero by Theorem~\ref{thm:charzero}.
These cases exhaust the possibilities. Each reduction preserves the
source multiplicity and preserves or decreases the target multiplicity,
or recombines the inequalities by additivity. Thus the inequality holds
for the original map.
\end{proof}

\begin{thebibliography}{99}

\bibitem{Artin}
M.~Artin,
\emph{Algebraic approximation of structures over complete local rings},
Publications Math\'ematiques de l'IH\'ES \textbf{36} (1969), 23--58.
\href{https://numdam.org/item/10.1007/BF02684596.pdf}{Original article}.

\bibitem{AFH}
L.~L.~Avramov, H.-B.~Foxby, and B.~Herzog,
\emph{Structure of local homomorphisms},
Journal of Algebra \textbf{164} (1994), 124--145.
\href{https://doi.org/10.1006/jabr.1994.1057}{doi:10.1006/jabr.1994.1057}.

\bibitem{BhattHochsterMa}
B.~Bhatt, M.~Hochster, and L.~Ma,
\emph{Lim Cohen--Macaulay sequences of modules},
\href{https://arxiv.org/abs/2410.18372v2}{arXiv:2410.18372v2}, 2026.
Numbered citations refer to this version.

\bibitem{BhattScholze}
B.~Bhatt and P.~Scholze,
\emph{Prisms and prismatic cohomology},
Annals of Mathematics (2) \textbf{196} (2022), no.~3, 1135--1275.
\href{https://doi.org/10.4007/annals.2022.196.3.5}{doi:10.4007/annals.2022.196.3.5}.
\href{https://arxiv.org/abs/1905.08229v4}{arXiv:1905.08229v4}.
Numbered citations refer to this version.

\bibitem{CLMST}
H.~Cai, S.~Lee, L.~Ma, K.~Schwede, and K.~Tucker,
\emph{Perfectoid signature, perfectoid Hilbert--Kunz multiplicity,
and an application to local fundamental groups},
Inventiones Mathematicae \textbf{244} (2026), 357--454.
\href{https://arxiv.org/abs/2209.04046v4}{arXiv:2209.04046v4}, 2025.
All numbered citations refer to this version.

\bibitem{Dutta}
S.~P.~Dutta,
\emph{Frobenius and multiplicities},
Journal of Algebra \textbf{85} (1983), no.~2, 424--448.
\href{https://doi.org/10.1016/0021-8693(83)90106-0}{doi:10.1016/0021-8693(83)90106-0}.

\bibitem{EisenbudErman}
D.~Eisenbud and D.~Erman,
\emph{Categorified duality in Boij--S\"oderberg theory and invariants
of free complexes},
Journal of the European Mathematical Society \textbf{19} (2017), 2657--2695.
\href{https://doi.org/10.4171/JEMS/725}{doi:10.4171/JEMS/725}.

\bibitem{EisenbudSchreyer}
D.~Eisenbud and F.-O.~Schreyer,
\emph{Betti numbers of graded modules and cohomology of vector bundles},
Journal of the American Mathematical Society \textbf{22} (2009), 859--888.
\href{https://arxiv.org/abs/0712.1843v3}{arXiv:0712.1843v3}, 2008.
Numbered citations refer to this version.

\bibitem{Faltings}
G.~Faltings,
\emph{Almost \'etale extensions},
Ast\'erisque \textbf{279} (2002), 185--270.
\href{https://numdam.org/item/AST_2002__279__185_0/}{Original article}.

\bibitem{Hanes}
D.~Hanes,
\emph{Special conditions on maximal Cohen--Macaulay modules, and
applications to the theory of multiplicities},
Ph.D. thesis, University of Michigan, 1999.

\bibitem{Heitmann}
R.~C.~Heitmann,
\emph{The plus closure in mixed characteristic},
Journal of Algebra \textbf{193} (1997), no.~2, 688--708.
\href{https://doi.org/10.1006/jabr.1997.7023}{doi:10.1006/jabr.1997.7023}.

\bibitem{HeitmannMa}
R.~C.~Heitmann and L.~Ma,
\emph{Extended plus closure in complete local rings},
Journal of Algebra \textbf{571} (2021), 134--150.
\href{https://arxiv.org/abs/1708.05761v3}{arXiv:1708.05761v3}, 2018.
All numbered citations refer to this version.

\bibitem{HunekeHK}
C.~Huneke,
\emph{Hilbert--Kunz multiplicity and the $F$-signature},
in \emph{Commutative Algebra}, I.~Peeva (ed.),
Springer, 2013, pp.~485--525.
\href{https://arxiv.org/abs/1409.0467v1}{arXiv:1409.0467v1}, 2014.
Numbered citations refer to this version.

\bibitem{IMW2022}
S.~B.~Iyengar, L.~Ma, and M.~E.~Walker,
\emph{Multiplicities and Betti numbers in local algebra via lim Ulrich points},
Algebra \& Number Theory \textbf{16} (2022), 1213--1257.
\href{https://doi.org/10.2140/ant.2022.16.1213}{doi:10.2140/ant.2022.16.1213}.
\href{https://arxiv.org/abs/2104.10140v2}{arXiv:2104.10140v2}.
Numbered citations refer to this version.

\bibitem{IMW2023}
S.~B.~Iyengar, L.~Ma, and M.~E.~Walker,
\emph{Lim Ulrich sequences and Boij--S\"oderberg cones},
Forum of Mathematics, Sigma \textbf{11} (2023), e115, 26 pp.
\href{https://doi.org/10.1017/fms.2023.108}{doi:10.1017/fms.2023.108}.

\bibitem{IMW2025}
S.~B.~Iyengar, L.~Ma, and M.~E.~Walker,
\emph{Ulrich modules over local rings of dimension two},
\href{https://arxiv.org/abs/2510.02698v1}{arXiv:2510.02698v1}, 2025.

\bibitem{KuranoShimomoto}
K.~Kurano and K.~Shimomoto,
\emph{An elementary proof of Cohen--Gabber theorem in the equal
characteristic $p>0$ case},
Tohoku Mathematical Journal \textbf{70} (2018), 377--389.
\href{https://doi.org/10.2748/tmj/1537495352}{doi:10.2748/tmj/1537495352}.

\bibitem{Lech}
C.~Lech,
\emph{Note on multiplicities of ideals},
Arkiv f\"or Matematik \textbf{4} (1960), 63--86.

\bibitem{Lech64}
C.~Lech,
\emph{Inequalities related to certain couples of local rings},
Acta Mathematica \textbf{112} (1964), 69--89.

\bibitem{Ma}
L.~Ma,
\emph{Lech's conjecture in dimension three},
Advances in Mathematics \textbf{322} (2017), 940--970.
\href{https://arxiv.org/abs/1609.00095v3}{arXiv:1609.00095v3}, 2018.
All numbered citations refer to this version.

\bibitem{MaUlrich}
L.~Ma,
\emph{Lim Ulrich sequences and Lech's conjecture},
Inventiones Mathematicae \textbf{231} (2023), 407--429.
\href{https://arxiv.org/abs/2005.02338v2}{arXiv:2005.02338v2}, 2022.
Numbered citations refer to this version.

\bibitem{Meng}
C.~Meng,
\emph{Strongly Lech-independent ideals and Lech's conjecture},
\href{https://arxiv.org/abs/2112.09849v3}{arXiv:2112.09849v3}, 2025.
Numbered citations refer to this version.

\bibitem{Monsky}
P.~Monsky,
\emph{The Hilbert--Kunz function},
Mathematische Annalen \textbf{263} (1983), 43--49.
\href{https://doi.org/10.1007/BF01457082}{doi:10.1007/BF01457082}.

\bibitem{SmallCM}
OpenAI,
\emph{A Complete Local Domain without a Small Cohen--Macaulay Module},
OpenAI Math Release preprint
\href{https://github.com/openai/math/blob/main/preprints/A-Complete-Local-Domain-Without-a-Small-Cohen-Macaulay-Module-September-23-2026/paper.pdf}{OAI:A-Complete-Local-Domain-Without-a-Small-Cohen-Macaulay-Module-September-23-2026},
2026.

\bibitem{PolstraTucker}
T.~Polstra and K.~Tucker,
\emph{$F$-signature and Hilbert--Kunz multiplicity:
a combined approach and comparison},
Algebra \& Number Theory \textbf{12} (2018), no.~1, 61--97.
\href{https://doi.org/10.2140/ant.2018.12.61}{doi:10.2140/ant.2018.12.61}.

\bibitem{Roberts}
P.~Roberts,
\emph{Intersection theorems},
in \emph{Commutative Algebra} (Berkeley, 1987),
Mathematical Sciences Research Institute Publications, vol.~15,
Springer, 1989, pp.~417--436.
\href{https://doi.org/10.1007/978-1-4612-3660-3_23}{doi:10.1007/978-1-4612-3660-3\_23}.

\bibitem{Shimomoto}
K.~Shimomoto,
\emph{The Frobenius action on local cohomology modules in mixed
characteristic},
Compositio Mathematica \textbf{143} (2007), no.~6, 1478--1492.
\href{https://doi.org/10.1112/S0010437X07002825}{doi:10.1112/S0010437X07002825}.

\bibitem{Stacks}
The Stacks Project Authors,
\emph{The Stacks Project}, 2026.
\href{https://stacks.math.columbia.edu}{stacks.math.columbia.edu}.

\bibitem{Yhee}
F.~C.~Yhee,
\emph{Ulrich modules and weakly lim Ulrich sequences do not always exist},
Forum of Mathematics, Sigma \textbf{11} (2023), e71.
\href{https://doi.org/10.1017/fms.2023.68}{doi:10.1017/fms.2023.68}.
Numbered citations refer to the published version.

\end{thebibliography}

\end{document}
